\documentclass[12pt]{article}
\usepackage[utf8]{inputenc}
\usepackage[spanish]{babel}
\usepackage{amsmath}
\usepackage{amssymb}
\usepackage{amsfonts}
\usepackage{tikz}
\usepackage{pgfplots}
\pgfplotsset{compat=1.18}
\usepackage{listings}
\usepackage{xcolor}

\lstset{
    language=Python,
    basicstyle=\ttfamily\small,
    keywordstyle=\color{blue},
    stringstyle=\color{red},
    showstringspaces=false,
    breaklines=true,
    frame=lines
}

\begin{document}

\subsection*{Enunciado}
¿Un objeto más pesado siempre tiene mayor rozamiento?

\subsection*{Solución}

\subsubsection*{Las Variables de la Fricción Seca}
En la física conceptual, es un error afirmar que el rozamiento depende de manera exclusiva y absoluta del peso del objeto. El modelo matemático de la fricción por deslizamiento (fricción de Coulomb) establece que esta fuerza disipativa ($f$) depende del producto de dos variables distintas: la Fuerza Normal ($N$) y el coeficiente de rozamiento ($\mu$). La ecuación es $f = \mu \cdot N$.

\subsubsection*{El Rol de la Normal frente al Peso}
La Fuerza Normal es la fuerza de compresión entre las dos superficies. Es cierto que, para un objeto apoyado en una superficie estrictamente horizontal, la magnitud de la Normal iguala a la del Peso. Sin embargo, si la superficie está inclinada, la Normal es solo una fracción del peso. Aún más drástico, si presionas un libro ligero contra una pared vertical, el peso apunta hacia abajo, pero la Normal (y por ende, la fricción que evita que caiga) depende de qué tan fuerte empujes horizontalmente, independientemente de cuán pesado sea el libro.

\subsubsection*{El Papel Fundamental de los Materiales}
El segundo factor, el coeficiente de rozamiento ($\mu$), cuantifica la interacción electromagnética y las rugosidades microscópicas entre los materiales específicos que entran en contacto. Un objeto sumamente pesado (como un bloque de hielo de 500 kg) deslizándose sobre una pista de patinaje (hielo sobre hielo, $\mu$ extremadamente bajo) puede experimentar una fuerza de rozamiento total mucho menor que un objeto ligero (como un neumático de caucho de 20 kg) arrastrado sobre asfalto seco ($\mu$ muy alto).

\subsubsection*{Visualización de la Paradoja del Peso y la Fricción}
Para refutar visualmente la afirmación del enunciado, el siguiente código en Python genera un gráfico de barras comparando dos escenarios hipotéticos que demuestran cómo un objeto más ligero puede generar más fricción si los materiales son altamente adherentes.

\begin{lstlisting}
import matplotlib.pyplot as plt

escenarios = ['Pesado sobre Hielo\n(N=1000 N, mu=0.05)', 'Ligero sobre Asfalto\n(N=200 N, mu=0.80)']
fuerzas_friccion = [1000 * 0.05, 200 * 0.80]

fig, ax = plt.subplots(figsize=(8, 5))
barras = ax.bar(escenarios, fuerzas_friccion, color=['lightblue', 'dimgray'], edgecolor='black', linewidth=1.5)

ax.set_ylabel('Fuerza de Rozamiento Resultante (Newtons)', fontsize=12, fontweight='bold')
ax.set_title('El Peso no Dicta la Friccion Absoluta', fontsize=14, fontweight='bold')
ax.set_ylim(0, 200)

for barra in barras:
    altura = barra.get_height()
    ax.text(barra.get_x() + barra.get_width()/2., altura + 5,
            f'f = {int(altura)} N',
            ha='center', va='bottom', fontsize=12, fontweight='bold')

ax.spines['top'].set_visible(False)
ax.spines['right'].set_visible(False)
ax.grid(axis='y', linestyle='--', alpha=0.5)

plt.tight_layout()
plt.show()
\end{lstlisting}

\subsubsection*{Conclusión}
No, un objeto más pesado no siempre tiene mayor rozamiento. La fuerza de rozamiento depende no solo de la fuerza de compresión (Fuerza Normal, que puede o no coincidir con el peso total), sino también críticamente del coeficiente de fricción entre los materiales en contacto.

\newpage

\subsection*{Enunciado}
¿Si no hay movimiento, no hay rozamiento?

\subsection*{Solución}

\subsubsection*{La Falsa Correlación entre Movimiento y Fricción}
Existe una concepción intuitiva errónea que asocia la fuerza de fricción exclusivamente con el acto de arrastrar o deslizar un objeto (fricción cinética). Sin embargo, las leyes de la mecánica dictaminan la existencia de una fuerza resistiva previa al movimiento, denominada fricción estática. 

\subsubsection*{La Dinámica de la Fricción Estática}
Imagina que empujas un refrigerador pesado y este no se mueve. Como tu empuje es una fuerza real aplicada en una dirección, y la aceleración es nula, la Primera Ley de Newton exige que exista una fuerza de igual magnitud pero en dirección opuesta que cancele tu empuje. Esa fuerza canceladora es el rozamiento estático. Esta interacción ocurre debido a que las irregularidades microscópicas de las superficies en contacto están "engranadas" y se resisten a ser cizalladas.

\subsubsection*{El Comportamiento Reactivo de la Superficie}
La fuerza de rozamiento estático posee una propiedad física fascinante: es reactiva y variable. Si aplicas 10 Newtons de fuerza y el objeto no se mueve, la fricción estática es exactamente de 10 Newtons. Si aumentas tu empuje a 50 Newtons y sigue sin moverse, la fricción estática crece automáticamente hasta 50 Newtons. Este crecimiento continúa hasta alcanzar un límite umbral de "agarre" ($f_{s,max} = \mu_s \cdot N$). Por lo tanto, no solo hay rozamiento cuando no hay movimiento, sino que suele ser la etapa donde el rozamiento alcanza su magnitud máxima.

\subsubsection*{Visualización de las Etapas de la Fricción}
Para comprender este comportamiento, el siguiente script en Python de Matplotlib grafica la evolución de la fuerza de rozamiento en función de la fuerza aplicada. En la gráfica se evidencia claramente una inmensa zona de rozamiento donde la velocidad es cero absoluto.

\begin{lstlisting}
import matplotlib.pyplot as plt
import numpy as np

fuerza_aplicada_estatica = np.linspace(0, 100, 100)
fuerza_rozamiento_estatica = fuerza_aplicada_estatica

fuerza_aplicada_cinetica = np.linspace(100, 150, 50)
fuerza_rozamiento_cinetica = np.full_like(fuerza_aplicada_cinetica, 75)

fig, ax = plt.subplots(figsize=(8, 5))

ax.plot(fuerza_aplicada_estatica, fuerza_rozamiento_estatica, color='red', linewidth=3, label='Friccion Estatica (v = 0)')
ax.plot(fuerza_aplicada_cinetica, fuerza_rozamiento_cinetica, color='blue', linewidth=3, label='Friccion Cinetica (v > 0)')
ax.plot([100, 100], [75, 100], color='gray', linestyle='--')

ax.text(40, 60, 'Friccion = Empuje\n(El objeto no se mueve)', fontsize=10, color='darkred', ha='center', rotation=45)
ax.plot(100, 100, 'ko')
ax.text(95, 105, 'Limite Estatico', fontsize=10, fontweight='bold')

ax.set_xlim(0, 150)
ax.set_ylim(0, 120)
ax.set_xlabel('Fuerza Externa Aplicada (N)', fontsize=12, fontweight='bold')
ax.set_ylabel('Fuerza de Rozamiento (N)', fontsize=12, fontweight='bold')
ax.set_title('Evolucion de la Friccion: Estatica vs Cinetica', fontsize=14, fontweight='bold')
ax.grid(True, linestyle='--', alpha=0.6)
ax.legend(loc='lower right')

plt.tight_layout()
plt.show()
\end{lstlisting}

\subsubsection*{Conclusión}
Falso. Puede existir una gran cantidad de rozamiento aunque no haya movimiento. A esta fuerza se le conoce como rozamiento estático, la cual se opone a cualquier fuerza externa aplicada, igualándola en magnitud para evitar que el cuerpo abandone su estado de reposo.

\newpage

\subsection*{Enunciado}
¿El aire no ejerce fuerza sobre objetos en movimiento?

\subsection*{Solución}

\subsubsection*{La Materialidad del Aire}
Un error fundamental en el pensamiento cotidiano es percibir el aire atmosférico como un "espacio vacío" o inmaterial. En la física conceptual, el aire es reconocido como un fluido real, compuesto por billones de moléculas de gases (principalmente nitrógeno y oxígeno) que poseen masa inercial y ocupan un volumen en el espacio.

\subsubsection*{El Principio de la Resistencia Aerodinámica}
Cuando un objeto sólido se mueve a través de la atmósfera, no se desplaza por la nada; debe abrirse paso chocando violentamente y apartando las moléculas de aire que se interponen en su trayectoria. De acuerdo con la Tercera Ley de Newton, por cada molécula de aire que el objeto empuja hacia adelante y hacia los lados (acción), la molécula de aire empuja al objeto hacia atrás (reacción). 
La sumatoria macroscópica de todos estos micro-impactos genera una fuerza mecánica masiva que se opone a la dirección del movimiento, conocida formalmente como resistencia del aire o arrastre aerodinámico (drag).

\subsubsection*{Variables de la Fuerza de Arrastre}
A diferencia de la fricción sólida que es relativamente independiente de la velocidad, la resistencia del aire es altamente dinámica. Crece exponencialmente a medida que el objeto aumenta su velocidad (choca contra más moléculas por segundo y con mayor ímpetu). También depende críticamente del área frontal del objeto (un paracaídas abierto abarca más moléculas que un paracaídas cerrado) y de su perfil aerodinámico. Es esta fuerza del aire la que impone un límite máximo de velocidad (velocidad terminal) a los paracaidistas en caída libre.

\subsubsection*{Visualización de la Fuerza del Aire en Dinámica}
Para ilustrar la importancia estructural de esta fuerza fluida, el siguiente script de Python genera un Diagrama de Cuerpo Libre de un objeto en caída a alta velocidad en el seno de la atmósfera.

\begin{lstlisting}
import matplotlib.pyplot as plt
import matplotlib.patches as patches

fig, ax = plt.subplots(figsize=(6, 7))

meteor = patches.Circle((0, 0), 0.8, color='darkorange', ec='black', lw=2)
ax.add_patch(meteor)
ax.text(0, 0, 'Objeto', ha='center', va='center', fontsize=12, fontweight='bold')

ax.annotate('', xy=(0, -3.5), xytext=(0, -0.8), 
            arrowprops=dict(facecolor='red', shrink=0, width=3, headwidth=10))
ax.text(0.3, -2.5, 'Gravedad (Peso)', color='red', fontsize=12, fontweight='bold')

ax.annotate('', xy=(0, 2.8), xytext=(0, 0.8), 
            arrowprops=dict(facecolor='blue', shrink=0, width=3, headwidth=10))
ax.text(0.3, 1.8, 'Resistencia\ndel Aire', color='blue', fontsize=12, fontweight='bold')

for x_line in [-1.2, -0.6, 0.6, 1.2]:
    ax.plot([x_line, x_line], [1, 3.5], color='skyblue', linewidth=1.5, linestyle='--')
    ax.plot([x_line, x_line], [-3.5, -1], color='skyblue', linewidth=1.5, linestyle='--')

ax.set_xlim(-4, 4)
ax.set_ylim(-4, 4)
ax.axis('off')
ax.set_title('El Aire como Fluido: Fuerza de Arrastre', fontsize=14, fontweight='bold')

plt.tight_layout()
plt.show()
\end{lstlisting}

\subsubsection*{Conclusión}
Falso. El aire ejerce una fuerza muy significativa sobre los objetos en movimiento. Al ser un fluido compuesto por moléculas con masa, el aire presenta una fricción mecánica (resistencia aerodinámica) contra cualquier cuerpo que intente desplazarlo, oponiéndose a su avance.

\newpage

\subsection*{Enunciado}
Una caja de 20 kg se encuentra en reposo sobre una superficie horizontal rugosa. Un hombre empuja horizontalmente la caja con una fuerza aplicada de 50 N hacia la derecha. Entre la superficie y la caja los coeficientes de rozamiento son: $\mu_e = 0.4$ y $\mu_c = 0.3$. Determine la fuerza de rozamiento que actúa sobre la caja.

\subsection*{Solución}

\subsubsection*{Comprendiendo los Coeficientes de Fricción}
En el estudio de la física conceptual, la fricción entre dos superficies sólidas se divide en dos regímenes: estático (cuando no hay deslizamiento relativo) y cinético (cuando las superficies se deslizan). El coeficiente de rozamiento estático ($\mu_e = 0.4$) dicta el límite máximo de resistencia antes de que el objeto rompa su estado de reposo. El coeficiente de rozamiento cinético ($\mu_c = 0.3$) dicta la resistencia constante una vez que el objeto ya está en movimiento. Es una regla general de la naturaleza que el agarre estático ($\mu_e$) siempre sea mayor que el roce dinámico ($\mu_c$).

\subsubsection*{Cálculo de la Fuerza Normal}
Para evaluar la fricción, primero debemos determinar qué tan fuertemente están presionadas las superficies entre sí. Al encontrarse en una superficie horizontal y sin fuerzas verticales externas, la caja se encuentra en equilibrio en el eje vertical ($\Sigma F_y = 0$). Esto significa que la Fuerza Normal ($N$) es exactamente igual al Peso ($P$) de la caja.
Utilizando la aceleración de la gravedad estándar $g = 9.8 \, \text{m/s}^2$:
\begin{equation}
    P = m \cdot g = 20 \, \text{kg} \cdot 9.8 \, \text{m/s}^2 = 196 \, \text{N}
\end{equation}
Por consiguiente, la Fuerza Normal es:
\begin{equation}
    N = 196 \, \text{N}
\end{equation}

\subsubsection*{Evaluación del Umbral de Fricción Estática}
Antes de asumir que la caja se mueve, debemos calcular el umbral máximo de resistencia que el suelo puede ofrecer. Este límite se calcula multiplicando el coeficiente estático por la normal:
\begin{equation}
    f_{e,\text{máx}} = \mu_e \cdot N
\end{equation}
\begin{equation}
    f_{e,\text{máx}} = 0.4 \cdot 196 \, \text{N}
\end{equation}
\begin{equation}
    f_{e,\text{máx}} = 78.4 \, \text{N}
\end{equation}
Esto significa que se requiere una fuerza externa estrictamente mayor a $78.4 \, \text{N}$ para lograr que la caja comience a deslizarse.

\subsubsection*{La Naturaleza Reactiva de la Fricción Estática}
El hombre aplica una fuerza de $50 \, \text{N}$. Al comparar esta fuerza impulsora con el umbral de resistencia ($50 \, \text{N} < 78.4 \, \text{N}$), es evidente que el empuje no es suficiente para romper el equilibrio. La caja permanecerá en reposo.
Aquí reside un concepto crucial: la fuerza de fricción estática no es de $78.4 \, \text{N}$. Si lo fuera, la caja aceleraría hacia atrás, lo cual es un absurdo físico. La fricción estática es una fuerza "inteligente" y reactiva; crece únicamente lo necesario para cancelar la fuerza aplicada, hasta su punto de ruptura. Dado que el hombre empuja con $50 \, \text{N}$, la superficie responde con exactamente $50 \, \text{N}$ en dirección opuesta para mantener la aceleración nula.

\subsubsection*{Visualización del Equilibrio Estático}
El siguiente código en Python utiliza Matplotlib para generar un Diagrama de Cuerpo Libre. En esta representación, es imperativo notar cómo el vector del empuje y el vector de la fricción estática son exactamente de la misma magnitud (se anulan mutuamente), evidenciando la Primera Ley de Newton.

\begin{lstlisting}
import matplotlib.pyplot as plt
import matplotlib.patches as patches

fig, ax = plt.subplots(figsize=(7, 5))

rect = patches.Rectangle((-1.5, -1), 3, 2, linewidth=1.5, edgecolor='black', facecolor='tan')
ax.add_patch(rect)
ax.text(0, 0, '20 kg', ha='center', va='center', fontsize=14, fontweight='bold')

ax.annotate('', xy=(4.0, 0), xytext=(1.5, 0), 
            arrowprops=dict(facecolor='green', shrink=0, width=2.5, headwidth=9))
ax.text(2.0, 0.3, 'Empuje (50 N)', color='green', fontsize=12, fontweight='bold')

ax.annotate('', xy=(-4.0, 0), xytext=(-1.5, 0), 
            arrowprops=dict(facecolor='orange', shrink=0, width=2.5, headwidth=9))
ax.text(-3.8, 0.3, 'Friccion (50 N)', color='orange', fontsize=12, fontweight='bold')

ax.annotate('', xy=(0, 3.5), xytext=(0, 1), 
            arrowprops=dict(facecolor='blue', shrink=0, width=2, headwidth=8))
ax.text(0.2, 2.5, 'Normal (196 N)', color='blue', fontsize=12, fontweight='bold')

ax.annotate('', xy=(0, -3.5), xytext=(0, -1), 
            arrowprops=dict(facecolor='red', shrink=0, width=2, headwidth=8))
ax.text(0.2, -2.8, 'Peso (196 N)', color='red', fontsize=12, fontweight='bold')

ax.set_xlim(-6, 6)
ax.set_ylim(-4.5, 4.5)
ax.axhline(0, color='gray', linestyle='--', linewidth=0.5, zorder=0)
ax.axvline(0, color='gray', linestyle='--', linewidth=0.5, zorder=0)
ax.set_title('DCL: Friccion Estatica Reactiva', fontsize=14, fontweight='bold')
ax.axis('off')

plt.tight_layout()
plt.show()
\end{lstlisting}

\subsubsection*{Conclusión}
La fuerza de rozamiento que actúa sobre la caja es de 50 N. Dado que la fuerza aplicada no supera el límite de fricción estática máxima (78.4 N), el objeto no se mueve, y la fuerza de fricción se ajusta para igualar exactamente y cancelar a la fuerza aplicada.

\newpage

\subsection*{Enunciado}
Un pequeño trozo de espuma de estireno (material de empaque) se suelta desde el reposo a una altura de 2 m sobre el suelo. Se sabe que, mientras acelera, la magnitud de su aceleración sigue el modelo $a = g - bv$. Tras haber caído una distancia de 0.5 m, el objeto alcanza su rapidez terminal ($v_T$) y, a partir de ese punto, tarda exactamente 5 s más en llegar al suelo. Determine el valor de la constante de resistencia $b$.

\subsection*{Solución}

\subsubsection*{Dinámica de la Resistencia del Aire}
En condiciones reales, cuando un objeto cae a través de la atmósfera, no experimenta una caída libre ideal. El aire actúa como un fluido que ejerce una fuerza de arrastre o resistencia aerodinámica en dirección opuesta al movimiento. Según la física conceptual, esta fuerza de arrastre incrementa conforme el objeto gana rapidez. 
El modelo matemático proporcionado ($a = g - bv$) describe precisamente este fenómeno: la aceleración neta ($a$) es igual a la aceleración gravitacional ($g$) menos un término de desaceleración ($bv$) que es directamente proporcional a la rapidez actual ($v$). La constante $b$ representa el factor de proporcionalidad que depende de la aerodinámica del objeto y la densidad del aire.

\subsubsection*{El Concepto Físico de Rapidez Terminal}
A medida que el trozo de espuma cae y su rapidez $v$ aumenta, el término $bv$ se hace cada vez más grande. Eventualmente, la resistencia del aire igualará exactamente a la fuerza gravitacional (el peso). Cuando esto ocurre, las fuerzas se equilibran, la fuerza neta se vuelve cero y, por la Segunda Ley de Newton, la aceleración también se vuelve cero ($a = 0$).
A esta rapidez máxima y constante se le denomina \textit{rapidez terminal} ($v_T$). Sustituyendo esta condición física en nuestro modelo:
\begin{equation}
    0 = g - b \cdot v_T
\end{equation}
Despejando analíticamente la constante $b$:
\begin{equation}
    b = \frac{g}{v_T}
\end{equation}

\subsubsection*{Análisis Cinemático de la Segunda Etapa}
Para encontrar el valor numérico de $b$, primero debemos determinar el valor de la rapidez terminal ($v_T$). El problema nos indica que el objeto alcanza $v_T$ tras caer $0.5 \, \text{m}$. Dado que la altura inicial total era de $2 \, \text{m}$, la distancia restante que el objeto recorre a rapidez terminal constante es:
\begin{equation}
    d_{\text{restante}} = 2 \, \text{m} - 0.5 \, \text{m} = 1.5 \, \text{m}
\end{equation}
El enunciado establece que recorre esta distancia restante en un tiempo exacto de $t = 5 \, \text{s}$. Al ser un movimiento con rapidez constante (aceleración nula), utilizamos la ecuación básica de la cinemática ($v = d/t$):
\begin{equation}
    v_T = \frac{1.5 \, \text{m}}{5 \, \text{s}} = 0.3 \, \text{m/s}
\end{equation}

\subsubsection*{Visualización de las Etapas de Caída}
Para comprender pedagógicamente el desglose del movimiento, el siguiente script de Python mediante Matplotlib ilustra las dos fases cinemáticas del objeto: la fase de aceleración variable y la fase de rapidez constante.

\begin{lstlisting}
import matplotlib.pyplot as plt
import matplotlib.patches as patches

fig, ax = plt.subplots(figsize=(6, 7))

ax.axhline(0, color='black', linewidth=2.5)
ax.text(0.5, -0.15, 'Suelo (0 m)', ha='center', fontsize=12, fontweight='bold')

ax.plot([-0.1, 0.1], [2, 2], color='black')
ax.text(-0.15, 2, '2.0 m', ha='right', va='center', fontsize=11, fontweight='bold')

ax.plot([-0.1, 0.1], [1.5, 1.5], color='black')
ax.text(-0.15, 1.5, '1.5 m', ha='right', va='center', fontsize=11, fontweight='bold')

ax.annotate('', xy=(0.5, 1.55), xytext=(0.5, 1.95), 
            arrowprops=dict(facecolor='red', shrink=0, width=2.5, headwidth=9))
ax.text(0.7, 1.75, 'Etapa 1: Aceleracion decreciente\nDistancia = 0.5 m\nv aumenta hasta v_T', color='red', va='center', fontsize=10, fontweight='bold')

ax.annotate('', xy=(0.5, 0.05), xytext=(0.5, 1.45), 
            arrowprops=dict(facecolor='blue', shrink=0, width=2.5, headwidth=9))
ax.text(0.7, 0.75, 'Etapa 2: Rapidez Terminal\nDistancia = 1.5 m\nv = 0.3 m/s (Constante)\nTiempo = 5 s', color='blue', va='center', fontsize=10, fontweight='bold')

ax.plot(0.5, 2.0, 's', color='silver', markersize=12, markeredgecolor='black')
ax.plot(0.5, 1.5, 's', color='silver', markersize=12, markeredgecolor='black')
ax.plot(0.5, 0.0, 's', color='silver', markersize=12, markeredgecolor='black')

ax.set_xlim(-0.5, 2.5)
ax.set_ylim(-0.3, 2.3)
ax.axis('off')
ax.set_title('Diagrama de Fases: Caida con Resistencia del Aire', fontsize=14, fontweight='bold')

plt.tight_layout()
plt.show()
\end{lstlisting}

\subsubsection*{Cálculo Cuantitativo de la Constante}
Habiendo determinado que la rapidez terminal es $v_T = 0.3 \, \text{m/s}$, procedemos a calcular la constante $b$. Utilizaremos la aceleración estándar de la gravedad en la superficie terrestre, $g = 9.8 \, \text{m/s}^2$.
Retomamos la ecuación despejada del equilibrio dinámico:
\begin{equation}
    b = \frac{g}{v_T}
\end{equation}
\begin{equation}
    b = \frac{9.8 \, \text{m/s}^2}{0.3 \, \text{m/s}}
\end{equation}
\begin{equation}
    b \approx 32.67 \, \text{s}^{-1}
\end{equation}
Las unidades resultantes son inversas de segundo ($1/\text{s}$ o $\text{s}^{-1}$), lo cual es físicamente coherente para que al multiplicar $b$ por la velocidad ($\text{m/s}$), el resultado final posea unidades de aceleración ($\text{m/s}^2$).

\subsubsection*{Conclusión}
El valor de la constante de resistencia aerodinámica $b$ para este trozo de espuma de estireno es de aproximadamente $32.67 \, \text{s}^{-1}$.

\newpage

\subsection*{Enunciado}
Usted intenta mover una caja pesada sobre un suelo rugoso. Si aplica una fuerza horizontal progresiva pero la caja aún no se mueve, ¿cómo se comporta la fuerza de fricción estática $\vec{fr}_e$?
\begin{enumerate}
    \item[a)] Es siempre menor que la fuerza aplicada, lo que permite que el objeto permanezca en reposo.
    \item[b)] Es igual a la mitad del peso del objeto multiplicado por el coeficiente de fricción.
    \item[c)] Desaparece una vez que la fuerza aplicada supera la fricción cinética.
    \item[d)] Aumenta en magnitud para igualar exactamente a la fuerza aplicada en todo momento.
    \item[e)] Se mantiene constante e igual a $\mu_e N$ desde el momento en que se empieza a empujar.
\end{enumerate}

\subsection*{Solución}

\subsubsection*{El Principio de Inercia y el Reposo}
De acuerdo con la Primera Ley de Newton, el estado de reposo de un objeto masivo solo puede preservarse si la suma vectorial de todas las fuerzas que actúan sobre él es exactamente cero ($\Sigma \vec{F} = 0$). Dado que el enunciado afirma explícitamente que se está aplicando una fuerza externa y la caja "aún no se mueve", la física conceptual exige la existencia simultánea de una fuerza en dirección opuesta que cancele exactamente el empuje aplicado. Esta fuerza es la fricción estática.

\subsubsection*{El Comportamiento Reactivo de la Fricción Estática}
Uno de los conceptos más malinterpretados en la dinámica es considerar a la fricción estática como una fuerza de magnitud constante. La fórmula $f_{e,\text{máx}} = \mu_e \cdot N$ no calcula la fuerza de fricción en todo momento, sino únicamente el umbral de ruptura (el límite máximo). 
Mientras no se alcance ese límite, la fricción estática actúa de manera "inteligente" o reactiva. Si aplicas 10 N, la fricción es de 10 N. Si aumentas progresivamente a 40 N, la fricción crece simultáneamente a 40 N. Debe igualar con exactitud meridiana a la fuerza aplicada; si fuera menor (opción a), habría una fuerza neta y el objeto aceleraría; si fuera constante a su valor máximo (opción e) desde el inicio, la caja aceleraría hacia atrás hacia la persona, lo cual es físicamente absurdo.

\subsubsection*{Visualización de la Relación Lineal}
Para demostrar pedagógicamente este comportamiento de ajuste perfecto, el siguiente código en Python modela gráficamente la magnitud de la fuerza de fricción en respuesta a una fuerza aplicada progresiva, ilustrando la pendiente de 1:1.

\begin{lstlisting}
import matplotlib.pyplot as plt
import numpy as np

fuerza_aplicada = np.linspace(0, 100, 100)
friccion_estatica = fuerza_aplicada 

fig, ax = plt.subplots(figsize=(6, 5))

ax.plot(fuerza_aplicada, friccion_estatica, color='darkorange', linewidth=3, label=r'$\vec{fr}_e$')

ax.set_xlim(0, 120)
ax.set_ylim(0, 120)
ax.plot([0, 120], [0, 120], 'k--', alpha=0.3) 

ax.set_xlabel('Fuerza Horizontal Aplicada (N)', fontsize=12, fontweight='bold')
ax.set_ylabel('Magnitud de Friccion Estatica (N)', fontsize=12, fontweight='bold')
ax.set_title('Comportamiento Reactivo de la Friccion', fontsize=14, fontweight='bold')

ax.text(45, 60, r'$fr_e = F_{aplicada}$', color='darkred', rotation=45, fontsize=12, fontweight='bold')
ax.plot(100, 100, 'ko', markersize=8)
ax.text(80, 105, 'Limite estatico ($\mu_e N$)', fontweight='bold')

ax.grid(True, linestyle='--', alpha=0.5)
plt.tight_layout()
plt.show()
\end{lstlisting}

\subsubsection*{Conclusión}
d) Aumenta en magnitud para igualar exactamente a la fuerza aplicada en todo momento.

\newpage

\subsection*{Enunciado}
Los bloques $A$ y $B$ de la figura se desplazan juntos con movimiento rectilíneo uniforme debido a la acción de la fuerza $\vec{F}$ (todas las superficies en contacto son rugosas). Entonces, necesariamente se cumple que:
\begin{enumerate}
    \item[a)] la $\vec{fr}_e$ que hace $A$ sobre $B$ está hacia la derecha.
    \item[b)] la $\vec{fr}_e$ que hace $A$ sobre $B$ está hacia la izquierda.
    \item[c)] la $\vec{fr}_e$ que hace $A$ sobre $B$ es nula.
    \item[d)] la $\vec{N}$ que hace $B$ sobre $A$ está hacia abajo.
    \item[e)] el bloque $A$ resbala hacia la izquierda del bloque $B$.
\end{enumerate}

\begin{center}
\begin{tikzpicture}[>=latex, scale=1.2]
    \draw[thick] (-1.5,0) -- (4,0);
    \draw[thick, fill=white] (0,0) rectangle (2,1);
    \node at (1,0.5) {$B$};
    \draw[thick, fill=white] (0.4,1) rectangle (1.6,1.8);
    \node at (1,1.4) {$A$};
    \draw[->, thick, line width=1.2pt] (2,0.5) -- (3.5,0.5) node[above] {$\vec{F}$};
\end{tikzpicture}
\end{center}

\subsection*{Solución}

\subsubsection*{Condiciones Cinemáticas del Sistema}
Para resolver el problema con rigor conceptual, la clave reside en decodificar la frase "se desplazan juntos con movimiento rectilíneo uniforme". El Movimiento Rectilíneo Uniforme (MRU) implica que la velocidad es constante en magnitud y dirección, lo cual cinemáticamente se traduce en una aceleración exactamente igual a cero ($\vec{a} = 0$).

\subsubsection*{Aislamiento y Análisis Dinámico del Bloque A}
Si evaluamos el sistema de forma aislada, comenzamos con el bloque superior $A$. Por la Segunda Ley de Newton, si su aceleración es cero, la sumatoria de fuerzas horizontales sobre él debe ser nula ($\Sigma F_{xA} = 0$). 
Dado que no existe ninguna cuerda, ni empuje externo actuando directamente sobre el bloque $A$, la única fuerza horizontal que podría existir es la fuerza de fricción estática entre la base de $A$ y el techo de $B$. Sin embargo, para que la suma de fuerzas sea cero y al no haber ninguna fuerza que contrarrestar, el valor de esta fuerza de fricción estática debe ser forzosamente nulo ($\vec{fr}_e = 0$). Si existiera cualquier fricción, el bloque $A$ aceleraría, violando la condición de MRU.

\subsubsection*{Aplicación del Principio de Acción y Reacción}
La Tercera Ley de Newton dicta que las fuerzas ocurren en pares interactivos. Si determinamos que el bloque $B$ ejerce una fuerza de fricción estática nula sobre el bloque $A$, entonces por par de acción y reacción, el bloque $A$ ejerce exactamente una fuerza de fricción estática nula sobre el bloque $B$.

\subsubsection*{Demostración Vectorial del Sistema}
El siguiente script en Python utiliza Matplotlib para separar los diagramas de cuerpo libre. Es fundamental visualizar que en el bloque $A$, la ausencia de vectores horizontales es el reflejo directo de su estado de equilibrio dinámico inercial.

\begin{lstlisting}
import matplotlib.pyplot as plt
import matplotlib.patches as patches

fig, (ax1, ax2) = plt.subplots(1, 2, figsize=(10, 5))

# DCL Bloque A
ax1.add_patch(patches.Rectangle((-1, -1), 2, 2, fill=True, color='lightblue', ec='black', lw=2))
ax1.text(0, 0, 'Bloque A\n(MRU, a = 0)', ha='center', va='center', fontweight='bold')
ax1.annotate('', xy=(0, 2.5), xytext=(0, 1), arrowprops=dict(facecolor='blue', shrink=0, width=2, headwidth=8))
ax1.text(0, 2.8, 'Normal (de B)', color='blue', ha='center', fontweight='bold')
ax1.annotate('', xy=(0, -2.5), xytext=(0, -1), arrowprops=dict(facecolor='red', shrink=0, width=2, headwidth=8))
ax1.text(0, -3.0, 'Peso (de A)', color='red', ha='center', fontweight='bold')
ax1.text(0, -4.0, 'Ninguna fuerza horizontal.\nFriccion estatica = 0', ha='center', fontweight='bold', color='darkred')
ax1.set_xlim(-4, 4); ax1.set_ylim(-4.5, 4.5); ax1.axis('off')

# DCL Bloque B
ax2.add_patch(patches.Rectangle((-1.5, -1), 3, 2, fill=True, color='tan', ec='black', lw=2))
ax2.text(0, 0, 'Bloque B\n(MRU, a = 0)', ha='center', va='center', fontweight='bold')
ax2.annotate('', xy=(3.5, 0), xytext=(1.5, 0), arrowprops=dict(facecolor='green', shrink=0, width=3, headwidth=10))
ax2.text(3.7, 0.4, '$\vec{F}$ Aplicada', color='green', ha='center', fontweight='bold')
ax2.annotate('', xy=(-3.5, 0), xytext=(-1.5, 0), arrowprops=dict(facecolor='orange', shrink=0, width=3, headwidth=10))
ax2.text(-3.7, 0.4, '$\vec{f}_{cinetica}$\n(del Suelo)', color='orange', ha='center', fontweight='bold')
ax2.text(0, -4.0, 'Fuerza Neta = 0\n$\vec{F}$ se cancela con $\vec{f}_{cinetica}$', ha='center', fontweight='bold')
ax2.set_xlim(-5, 5); ax2.set_ylim(-4.5, 4.5); ax2.axis('off')

plt.tight_layout()
plt.show()
\end{lstlisting}

\subsubsection*{Conclusión}
c) la $\vec{fr}_e$ que hace $A$ sobre $B$ es nula. Al desplazarse a velocidad constante, la aceleración de A es cero, requiriendo que la fuerza neta horizontal sea nula. Como la fricción es la única interacción horizontal posible para A, su valor debe ser cero, anulando su correspondiente par de reacción sobre B.

\newpage

\subsection*{Enunciado}
En la figura: $F = 20 \, \text{N}$, $m = 1 \, \text{kg}$ y $\mu = 0.6$. La fuerza de rozamiento que la superficie rugosa ejerce sobre el bloque es:
\begin{enumerate}
    \item[a)] cinética, de $12 \, \text{N}$
    \item[b)] cinética, de $4 \, \text{N}$
    \item[c)] estática, de $10 \, \text{N}$
    \item[d)] estática, de $12 \, \text{N}$
    \item[e)] nula
\end{enumerate}

\begin{center}
\begin{tikzpicture}[>=latex, scale=1.2]
    \draw[thick] (0,-1.5) -- (0,1.5);
    \draw[thick, fill=white] (0,-0.5) rectangle (1,0.5);
    \node at (0.5,0) {$m$};
    \draw[<-, thick] (1,0) -- (2.5,0) node[above, midway] {$\vec{F}$};
    \draw[->, thick, bend left=45] (0.2,0.6) to (0.8,0.9);
    \node at (1.0, 1.0) {$\mu$};
\end{tikzpicture}
\end{center}

\subsection*{Solución}

\subsubsection*{El Eje Horizontal y la Fuerza Normal}
Para resolver este sistema, debemos abandonar la noción preconcebida de que la Fuerza Normal siempre es igual al peso. La Fuerza Normal es estrictamente una fuerza de reacción perpendicular a la superficie de contacto. En este escenario, la superficie es una pared vertical. 
La fuerza externa $\vec{F}$ empuja el bloque horizontalmente contra la pared con una magnitud de $20 \, \text{N}$. Como el bloque no atraviesa la pared ni rebota alejándose de ella, se encuentra en equilibrio en el eje horizontal ($\Sigma F_x = 0$). Por la Tercera Ley de Newton, la pared debe reaccionar empujando el bloque en sentido contrario con una Fuerza Normal de magnitud idéntica:
\begin{equation}
    N = F = 20 \, \text{N}
\end{equation}

\subsubsection*{El Eje Vertical y el Límite de Fricción}
En el eje vertical, la gravedad tira del bloque hacia abajo. Asumiendo el valor estándar simplificado de $g = 10 \, \text{m/s}^2$ (común en este tipo de opciones múltiples para obtener valores enteros), el peso del bloque es:
\begin{equation}
    P = m \cdot g = 1 \, \text{kg} \cdot 10 \, \text{m/s}^2 = 10 \, \text{N}
\end{equation}
La fuerza que evita que el bloque caiga es la fricción estática, la cual apunta hacia arriba paralela a la pared. Antes de determinar su valor real, debemos calcular su capacidad máxima de retención (el umbral estático):
\begin{equation}
    f_{e,\text{máx}} = \mu \cdot N
\end{equation}
\begin{equation}
    f_{e,\text{máx}} = 0.6 \cdot 20 \, \text{N} = 12 \, \text{N}
\end{equation}
La pared es capaz de soportar hasta $12 \, \text{N}$ de peso sin permitir el deslizamiento.

\subsubsection*{La Naturaleza Reactiva del Rozamiento Estático}
Comparamos la fuerza que intenta mover el bloque (el peso de $10 \, \text{N}$) con la resistencia máxima disponible ($12 \, \text{N}$). Dado que el peso es menor que la capacidad máxima de fricción ($10 \, \text{N} < 12 \, \text{N}$), el bloque no se desliza. Por lo tanto, el estado del rozamiento es indudablemente \textbf{estático}.
Sin embargo, un error conceptual gravísimo sería afirmar que la fricción estática es de $12 \, \text{N}$. Si lo fuera, el bloque tendría una fuerza neta de $2 \, \text{N}$ hacia arriba y comenzaría a levitar desafiando la gravedad. La fricción estática es una fuerza reactiva que crece solo lo necesario para equilibrar las fuerzas. Si el peso tira hacia abajo con $10 \, \text{N}$, la fricción estática se ajusta a exactamente $10 \, \text{N}$ hacia arriba para mantener el reposo ($\Sigma F_y = 0$).

\subsubsection*{Visualización de Fuerzas Ortogonales}
El siguiente código en Python genera el Diagrama de Cuerpo Libre utilizando Matplotlib, ilustrando cómo el eje de compresión (horizontal) dicta el límite del eje de movimiento inminente (vertical).

\begin{lstlisting}
import matplotlib.pyplot as plt
import matplotlib.patches as patches

fig, ax = plt.subplots(figsize=(6, 6))

bloque = patches.Rectangle((-1, -1), 2, 2, fill=True, color='lightskyblue', ec='black', lw=2)
ax.add_patch(bloque)
ax.text(0, 0, '1 kg', ha='center', va='center', fontweight='bold', fontsize=14)

ax.annotate('', xy=(-3.5, 0), xytext=(-1.1, 0), 
            arrowprops=dict(facecolor='blue', shrink=0, width=2, headwidth=8))
ax.text(-3.0, 0.3, 'Normal (20 N)', color='blue', fontweight='bold')

ax.annotate('', xy=(-1.1, 0), xytext=(-3.5, 0), 
            arrowprops=dict(facecolor='green', shrink=0, width=2, headwidth=8))
ax.text(-2.5, -0.5, 'Empuje (20 N)', color='green', fontweight='bold')

ax.annotate('', xy=(0, 2.8), xytext=(0, 1.1), 
            arrowprops=dict(facecolor='orange', shrink=0, width=2, headwidth=8))
ax.text(0.2, 2.0, 'Friccion Estatica (10 N)', color='orange', fontweight='bold')

ax.annotate('', xy=(0, -2.8), xytext=(0, -1.1), 
            arrowprops=dict(facecolor='red', shrink=0, width=2, headwidth=8))
ax.text(0.2, -2.0, 'Peso (10 N)', color='red', fontweight='bold')

ax.plot([1.2, 1.2], [-3, 3], color='black', linewidth=4)
ax.text(1.5, 0, 'Pared', rotation=-90, fontweight='bold', fontsize=12, va='center')

ax.set_xlim(-4.5, 2.5)
ax.set_ylim(-3.5, 3.5)
ax.axis('off')
ax.set_title('DCL: Friccion Estatica en Pared Vertical', fontsize=14, fontweight='bold')

plt.tight_layout()
plt.show()
\end{lstlisting}

\subsubsection*{Conclusión}
c) estática, de $10 \, \text{N}$. Al no superar la fuerza máxima de adherencia ($12 \, \text{N}$), el bloque permanece en reposo (fricción estática) y la magnitud de esta fricción se ajusta exactamente para contrarrestar el peso ($10 \, \text{N}$).

\newpage

\subsection*{Enunciado}
Un objeto de gran tamaño cae a través del aire y alcanza su rapidez terminal. ¿Cuál de las siguientes afirmaciones describe correctamente el estado de las fuerzas en ese momento?
\begin{enumerate}
    \item[a)] La fuerza de resistencia del aire desaparece debido a la alta velocidad.
    \item[b)] La magnitud de la fuerza resistiva del aire es igual al peso del objeto.
    \item[c)] El objeto se detiene instantáneamente en el aire por la fricción del fluido.
    \item[d)] La fuerza neta sobre el objeto es máxima y apunta hacia arriba.
    \item[e)] La aceleración del objeto es igual a la gravedad.
\end{enumerate}

\subsection*{Solución}

\subsubsection*{Dinámica de la Caída en un Fluido}
Para evaluar esta situación desde los principios de Newton, debemos identificar las fuerzas en conflicto. Al caer a través de la atmósfera, un objeto masivo experimenta su Peso (fuerza de atracción gravitacional) empujándolo constantemente hacia abajo. Simultáneamente, el fluido atmosférico presenta una Resistencia del Aire (arrastre aerodinámico) que empuja hacia arriba, oponiéndose al descenso.

\subsubsection*{El Concepto de Rapidez Terminal}
La resistencia del aire es una fuerza dinámica; es decir, no es constante. Su magnitud depende críticamente de la rapidez del objeto. En el instante exacto en que el objeto es soltado, su rapidez es cero, por lo que la resistencia del aire es nula y la aceleración es máxima (igual a la gravedad $g$). 
Sin embargo, conforme el objeto acelera y gana rapidez, comienza a chocar contra más moléculas de aire por segundo y con mayor violencia. Esto causa que la fuerza de resistencia del aire crezca progresivamente.

\subsubsection*{Análisis del Equilibrio Dinámico}
Llegará un momento físico inevitable en el cual la fuerza de resistencia del aire orientada hacia arriba crecerá tanto que igualará exactamente en magnitud al peso del objeto orientado hacia abajo.
\begin{equation}
    f_{\text{aire}} = P
\end{equation}
Al plantear la sumatoria de fuerzas verticales, encontramos que se anulan por completo:
\begin{equation}
    \Sigma F_y = f_{\text{aire}} - P = 0
\end{equation}
De acuerdo con la Segunda Ley de Newton, si la fuerza neta es cero, la aceleración se vuelve instantáneamente cero ($\vec{a} = 0$). Sin aceleración, el objeto ya no puede seguir ganando rapidez. Su velocidad se "estanca" en su valor máximo alcanzable, estado que en física se denomina formalmente \textit{rapidez terminal}. El objeto no se detiene; simplemente deja de acelerar, continuando su caída hacia el suelo a una rapidez alta y constante.

\subsubsection*{Conclusión}
b) La magnitud de la fuerza resistiva del aire es igual al peso del objeto. Al alcanzar la rapidez terminal, las fuerzas opuestas se equilibran (fuerza neta nula), provocando que la aceleración sea cero y el objeto mantenga una velocidad constante en su descenso final.

\newpage

\subsection*{Enunciado}
El modelo de resistencia de fluidos donde la fuerza es proporcional al cuadrado de la rapidez ($R = \frac{1}{2}D\rho Av^2$) es más adecuado para:
\begin{enumerate}
    \item[a)] partículas de polvo moviéndose lentamente.
    \item[b)] objetos grandes a altas velocidades, como un paracaidista.
    \item[c)] una esfera pequeña en un líquido muy viscoso como miel.
    \item[d)] un objeto en el vacío.
    \item[e)] cualquier situación que implique el aire.
\end{enumerate}

\subsection*{Solución}

\subsubsection*{Los Dos Regímenes de la Resistencia de Fluidos}
Cuando un objeto se mueve a través de un fluido (ya sea un líquido o un gas como la atmósfera terrestre), experimenta una fuerza de arrastre que se opone a su movimiento. En la física conceptual y la mecánica de fluidos, esta resistencia disipativa generalmente se clasifica en dos modelos matemáticos distintos. Qué modelo usar depende de cuán dominante sea la inercia del fluido frente a la viscosidad del mismo.

\subsubsection*{Modelo de Arrastre Viscoso (Proporcional a $v$)}
Para objetos microscópicos (o muy pequeños) que se mueven a velocidades extremadamente bajas, o para objetos desplazándose en fluidos sumamente espesos, las fuerzas de fricción de la viscosidad son las que dominan el entorno. En este régimen analítico (flujo laminar), el fluido se desliza ordenadamente alrededor del objeto sin dejar turbulencias bruscas. 
La resistencia aerodinámica se modela matemáticamente como directamente proporcional a la rapidez ($R \propto v$ o $R = bv$, conocido como la Ley de Stokes). Esto descarta las opciones (a) y (c).

\subsubsection*{Modelo de Arrastre Inercial (Proporcional a $v^2$)}
Para objetos de tamaño macroscópico (escala humana) que se desplazan a velocidades moderadas o altas en fluidos de baja viscosidad como el aire o el agua, la inercia de la masa del fluido es el factor dominante. El objeto avanza apartando violentamente a las moléculas del fluido y genera una estela turbulenta y caótica detrás de él. 
En este escenario inercial, el modelo empírico y preciso es la ecuación de arrastre de Newton: $R = \frac{1}{2}D\rho A v^2$. Esta fuerza crece drásticamente en proporción al cuadrado de la rapidez ($v^2$), al área frontal geométrica expuesta ($A$) y a la densidad del medio ($\rho$). Este es exactamente el comportamiento dinámico de un paracaidista en caída libre, un automóvil desplazándose por una autopista o un avión comercial en vuelo.

\subsubsection*{Visualización de los Modelos de Arrastre}
Para ilustrar la marcada diferencia en cómo escala la fuerza resistiva a medida que el objeto acelera, el siguiente script en Python utiliza Matplotlib para contrastar el crecimiento lineal (Stokes) contra el crecimiento exponencial cuadrático (Newton).

\begin{lstlisting}
import matplotlib.pyplot as plt
import numpy as np

v = np.linspace(0, 10, 100)
arrastre_viscoso = 2 * v          
arrastre_inercial = 0.5 * v**2    

fig, ax = plt.subplots(figsize=(7, 4.5))

ax.plot(v, arrastre_viscoso, color='green', linewidth=2.5, label=r'Arrastre Viscoso ($R \propto v$)')
ax.plot(v, arrastre_inercial, color='red', linewidth=2.5, label=r'Arrastre Inercial ($R \propto v^2$)')

ax.set_xlim(0, 10)
ax.set_ylim(0, 20)
ax.set_xlabel('Rapidez (v)', fontsize=12, fontweight='bold')
ax.set_ylabel('Fuerza de Resistencia (R)', fontsize=12, fontweight='bold')
ax.set_title('Comparativa: Modelos de Friccion de Fluidos', fontsize=14, fontweight='bold')

ax.grid(True, linestyle='--', alpha=0.6)
ax.legend(loc='upper left', fontsize=11)

plt.tight_layout()
plt.show()
\end{lstlisting}

\subsubsection*{Evaluación de Casos Atípicos}
Por rigor técnico, evaluamos las opciones restantes. La opción (d) es físicamente una imposibilidad absoluta porque en el vacío puro no existe gas material (la densidad es $\rho = 0$) que ofrezca ningún tipo de resistencia, por lo que $R=0$. La opción (e) es una sobregeneralización falsa, ya que una mota de polvo flotando en el aire obedece al modelo lineal de Stokes, no al modelo cuadrático, a pesar de estar rodeada de aire.

\subsubsection*{Conclusión}
b) objetos grandes a altas velocidades, como un paracaidista. El modelo cuadrático de $v^2$ cuantifica con exactitud matemática la resistencia inercial y turbulenta característica de los objetos de gran envergadura (escala macroscópica) desplazándose velozmente a través de fluidos.

\newpage

\subsection*{Enunciado}
Una caja de 20 kg se encuentra en reposo sobre una superficie horizontal rugosa. Un hombre empuja horizontalmente la caja con una fuerza $\vec{F}$ de 50 N dirigida hacia la derecha. Entre la superficie y la caja, los coeficientes de rozamiento son: $\mu_e = 0.4$ y $\mu_c = 0.3$. Determine la fuerza de rozamiento que actúa sobre la caja.

\subsection*{Solución}

\subsubsection*{Determinación de la Fuerza Normal Transversal}
Para iniciar el análisis friccional, el primer pilar es determinar la fuerza de soporte o compresión perpendicular entre la caja y el suelo rugoso. Dado que la caja reposa en equilibrio mecánico en el eje vertical ($\Sigma F_y = 0$), y la superficie es perfectamente horizontal, la magnitud de la Fuerza Normal ($N$) contrarresta exactamente el Peso gravitacional ($P$).
Utilizando la aceleración estándar de $g = 9.8 \, \text{m/s}^2$:
\begin{equation}
    N = P = m \cdot g = 20 \, \text{kg} \cdot 9.8 \, \text{m/s}^2 = 196 \, \text{N}
\end{equation}

\subsubsection*{Cálculo Cuantitativo del Límite de Fricción Estática}
La mecánica clásica estipula que un objeto requiere vencer un umbral específico de adherencia microscópica antes de que el movimiento relativo comience. Cuantificamos esta resistencia máxima estática utilizando su coeficiente empírico correspondiente ($\mu_e$):
\begin{equation}
    f_{e,\text{máx}} = \mu_e \cdot N
\end{equation}
\begin{equation}
    f_{e,\text{máx}} = 0.4 \cdot 196 \, \text{N} = 78.4 \, \text{N}
\end{equation}

\subsubsection*{Comprobación del Estado Dinámico Inercial}
La fuerza física real que empuja el bloque hacia la derecha es de $50 \, \text{N}$. Al contrastar esta fuerza impulsora contra la capacidad máxima de resistencia estructural ($50 \, \text{N} < 78.4 \, \text{N}$), la matemática de fuerzas revela que la energía aplicada es insuficiente para quebrar el agarre del suelo. El bloque preserva su estado de reposo absoluto. 

\subsubsection*{La Regla Reactiva Vectorial de la Fricción}
Concluido el estado, se aplica la máxima de la Primera Ley de Newton: para que exista reposo (aceleración nula), la sumatoria de fuerzas horizontales debe ser matemáticamente cero. Esto certifica que la fricción estática no está operando a su límite de $78.4 \, \text{N}$ (lo cual empujaría la caja hacia atrás), sino que actúa de forma gradualmente reactiva. Crece lo estrictamente indispensable para igualar a la fuerza impulsora. 
Adicionalmente, debido a la Tercera Ley de Newton, el vector de fricción reacciona en dirección geométricamente opuesta al estímulo; si el empuje transcurre hacia la derecha, el suelo resiste en vector hacia la izquierda.

\subsubsection*{Conclusión}
La fuerza de rozamiento que actúa disipativamente sobre la caja es de $50 \, \text{N}$, orientada hacia la izquierda.

\newpage

\subsection*{Enunciado}
Una fuerza horizontal $\vec{F}$ es aplicada como se muestra en la figura. El bloque de abajo tiene una masa $m_B = 16 \, \text{kg}$ y descansa sobre una mesa lisa. El bloque de arriba tiene una masa $m_A = 8 \, \text{kg}$. Suponga que el coeficiente de fricción estático entre los dos bloques es $\mu_e = 0.2$. Use $g = 10 \, \text{m/s}^2$. Determine:
\begin{enumerate}
    \item[a)] la máxima magnitud de $\vec{F}$, de tal manera que los dos bloques se muevan juntos sin deslizarse, y
    \item[b)] la magnitud de la aceleración correspondiente.
\end{enumerate}

\begin{center}
\begin{tikzpicture}[>=latex, scale=1.0]
    \draw[thick, fill=red!20] (-2,0) rectangle (4,-0.2);
    \draw[thick, fill=cyan!60!black, text=white] (0,0) rectangle (3,1.5);
    \node at (1.5, 0.75) {\textbf{$B$}};
    \draw[thick, fill=blue!20] (0.5,1.5) rectangle (2.5,2.5);
    \node at (1.5, 2.0) {\textbf{$A$}};
    \draw[->, thick, line width=1.5pt] (-1.5,0.75) -- (0,0.75) node[above, midway] {$\vec{F}$};
\end{tikzpicture}
\end{center}

\subsection*{Solución}

\subsubsection*{El Rol Motor de la Fricción (Bloque Superior)}
En la física conceptual, el primer paso es identificar qué fuerza causa el movimiento de cada masa. La fuerza externa $\vec{F}$ se aplica directa y exclusivamente sobre el bloque inferior $B$. El bloque superior $A$ no está siendo empujado por ninguna mano ni cuerda. Si la superficie entre ambos bloques fuera perfectamente lisa (hielo sobre hielo), al empujar $B$, el bloque $A$ se quedaría en su lugar por inercia y $B$ resbalaría por debajo de él. 
Por lo tanto, la única fuerza física capaz de arrastrar al bloque $A$ hacia adelante y obligarlo a acelerar junto con $B$ es la fuerza de \textbf{fricción estática} generada en la superficie de contacto entre ambos. 

\subsubsection*{El Límite de Adherencia y Aceleración Máxima}
Para que los bloques "se muevan juntos sin deslizarse", la fricción entre ellos debe mantenerse en el régimen estático. Sin embargo, esta fuerza de agarre tiene un límite físico máximo ($f_{e,\text{máx}} = \mu_e \cdot N_A$).
Primero, calculamos la Fuerza Normal que soporta al bloque $A$, la cual equilibra su peso:
\begin{equation}
    N_A = m_A \cdot g = 8 \, \text{kg} \cdot 10 \, \text{m/s}^2 = 80 \, \text{N}
\end{equation}
Ahora, calculamos la tracción máxima (fuerza de fricción estática máxima) disponible antes de que $A$ resbale:
\begin{equation}
    f_{e,\text{máx}} = 0.2 \cdot 80 \, \text{N} = 16 \, \text{N}
\end{equation}
Esta es la fuerza máxima absoluta que puede empujar al bloque $A$ hacia adelante. Aplicando la Segunda Ley de Newton solo para el bloque $A$, descubrimos su aceleración límite:
\begin{equation}
    a_{\text{máx}} = \frac{f_{e,\text{máx}}}{m_A}
\end{equation}
\begin{equation}
    a_{\text{máx}} = \frac{16 \, \text{N}}{8 \, \text{kg}} = 2 \, \text{m/s}^2
\end{equation}
Si el sistema acelera a más de $2 \, \text{m/s}^2$, el bloque $A$ no tendrá suficiente fricción para mantenerse al ritmo y resbalará hacia atrás respecto a $B$.

\subsubsection*{El Sistema Combinado y la Fuerza Total}
Una vez que conocemos la aceleración máxima permitida ($2 \, \text{m/s}^2$), podemos tratar a ambos bloques como si estuvieran pegados, formando un solo gran bloque o "sistema combinado", puesto que se mueven juntos en este límite. 
La masa total del sistema inercial es:
\begin{equation}
    M_{\text{total}} = m_A + m_B = 8 \, \text{kg} + 16 \, \text{kg} = 24 \, \text{kg}
\end{equation}
La única fuerza externa horizontal que mueve a este sistema completo sobre la mesa lisa es $\vec{F}$. Aplicando nuevamente la Segunda Ley de Newton al sistema global:
\begin{equation}
    F = M_{\text{total}} \cdot a_{\text{máx}}
\end{equation}
\begin{equation}
    F = 24 \, \text{kg} \cdot 2 \, \text{m/s}^2
\end{equation}
\begin{equation}
    F = 48 \, \text{N}
\end{equation}

\subsubsection*{Visualización de las Interacciones Internas}
Para consolidar este principio de acción y reacción interno, el siguiente script en Python mediante Matplotlib separa el sistema en un Diagrama de Cuerpo Libre por componentes, mostrando cómo la fricción actúa como par interactivo.

\begin{lstlisting}
import matplotlib.pyplot as plt
import matplotlib.patches as patches

fig, (ax1, ax2) = plt.subplots(1, 2, figsize=(10, 4.5))

# DCL Bloque A
ax1.add_patch(patches.Rectangle((-1, -1), 2, 2, fill=True, color='lavender', ec='black', lw=2))
ax1.text(0, 0, 'Bloque A\n(8 kg)', ha='center', va='center', fontweight='bold')
ax1.annotate('', xy=(3, 0), xytext=(1, 0), arrowprops=dict(facecolor='orange', shrink=0, width=2, headwidth=8))
ax1.text(2, 0.4, 'Friccion (16 N)\n"Fuerza Motriz"', color='darkorange', ha='center', fontweight='bold', fontsize=10)
ax1.set_xlim(-4, 4); ax1.set_ylim(-3, 3); ax1.axis('off')
ax1.set_title('Bloque A: Arrastrado por la Friccion', fontsize=12, fontweight='bold')

# DCL Bloque B
ax2.add_patch(patches.Rectangle((-1.5, -1), 3, 2, fill=True, color='steelblue', ec='black', lw=2))
ax2.text(0, 0, 'Bloque B\n(16 kg)', ha='center', va='center', fontweight='bold', color='white')
ax2.annotate('', xy=(-3.5, 0), xytext=(-1.5, 0), arrowprops=dict(facecolor='green', shrink=0, width=3, headwidth=10))
ax2.text(-2.5, 0.4, 'Empuje Externa (48 N)', color='green', ha='center', fontweight='bold', fontsize=10)
ax2.annotate('', xy=(-3.5, -0.6), xytext=(-1.5, -0.6), arrowprops=dict(facecolor='red', shrink=0, width=2, headwidth=8))
ax2.text(-2.5, -1.2, 'Friccion Reaccion (-16 N)', color='red', ha='center', fontweight='bold', fontsize=10)
ax2.set_xlim(-5, 4); ax2.set_ylim(-3, 3); ax2.axis('off')
ax2.set_title('Bloque B: Empujado y Frenado', fontsize=12, fontweight='bold')

plt.tight_layout()
plt.show()
\end{lstlisting}

\subsubsection*{Conclusión}
\begin{enumerate}
    \item[a)] La máxima magnitud de la fuerza $\vec{F}$ para que no haya deslizamiento es de $48 \, \text{N}$.
    \item[b)] La magnitud de la aceleración correspondiente del sistema en ese punto límite es de $2 \, \text{m/s}^2$.
\end{enumerate}

\newpage

\subsection*{Enunciado}
El bloque $A$ de la figura tiene una masa de $2 \, \text{kg}$ y el bloque $B$ tiene una masa de $4 \, \text{kg}$. El coeficiente de fricción entre todas las superficies es de $0.3$. Determine la magnitud de la fuerza horizontal $\vec{F}$ necesaria para arrastrar a $B$ hacia la izquierda con rapidez constante. Asuma que $A$ y $B$ están conectados mediante una cuerda ideal apoyada sobre una polea ideal.

\subsection*{Solución}

\subsubsection*{El Principio de Aislamiento y la Cinemática}
Para resolver sistemas compuestos por múltiples objetos interactuando, el pilar de la dinámica newtoniana es el diagrama de cuerpo libre: debemos analizar cada bloque por separado. 
El enunciado nos da una clave cinemática fundamental: el sistema se mueve con "rapidez constante". Esto significa que ninguno de los bloques está acelerando ($a = 0$). Por lo tanto, de acuerdo con la Primera Ley de Newton, la fuerza neta en ambos bloques debe ser exactamente cero tanto en el eje vertical como en el horizontal.

\subsubsection*{Análisis Dinámico del Bloque A (El bloque superior)}
Al observar el bloque $A$, notamos que mientras $B$ es arrastrado hacia la izquierda, la polea y la cuerda obligan al bloque $A$ a moverse hacia la derecha (en relación con el suelo).

En el eje vertical, el bloque $A$ no levita ni se hunde. La gravedad tira de él hacia abajo con su peso ($P_A$), y el bloque $B$ lo sostiene empujándolo hacia arriba con una Fuerza Normal ($N_A$). Utilizando $g = 10 \, \text{m/s}^2$:
\begin{equation}
    N_A = m_A \cdot g = 2 \, \text{kg} \cdot 10 \, \text{m/s}^2 = 20 \, \text{N}
\end{equation}

En el eje horizontal, como el bloque $A$ se desliza hacia la derecha, la fricción cinética ($f_A$) ejercida por el bloque $B$ actúa hacia la izquierda para oponerse a ese movimiento. Su magnitud es:
\begin{equation}
    f_A = \mu \cdot N_A = 0.3 \cdot 20 \, \text{N} = 6 \, \text{N}
\end{equation}

Para que el bloque $A$ se mueva a rapidez constante, la Tensión ($T$) de la cuerda que lo jala hacia la derecha debe equilibrar exactamente a esta fricción hacia la izquierda:
\begin{equation}
    T = f_A = 6 \, \text{N}
\end{equation}

\subsubsection*{Análisis Dinámico del Bloque B (El bloque inferior)}
El bloque $B$ se mueve hacia la izquierda. Debemos identificar rigurosamente todas las fuerzas que interactúan con él.

En el eje vertical, el bloque $B$ soporta su propio peso ($P_B = 40 \, \text{N}$) más la fuerza con la que el bloque $A$ presiona sobre él hacia abajo ($20 \, \text{N}$). Por lo tanto, el suelo debe sostener a ambos, ejerciendo una Fuerza Normal ($N_B$) equivalente a la suma de ambos pesos:
\begin{equation}
    N_B = 40 \, \text{N} + 20 \, \text{N} = 60 \, \text{N}
\end{equation}

En el eje horizontal, el bloque $B$ es jalado hacia la izquierda por la fuerza externa $\vec{F}$. Existen tres fuerzas distintas que se oponen a esto apuntando hacia la derecha:
\begin{itemize}
    \item \textbf{La Tensión de la cuerda ($T$):} La cuerda jala al bloque $B$ hacia la derecha. Como es una polea ideal, la tensión es la misma en toda la cuerda: $T = 6 \, \text{N}$.
    \item \textbf{La fricción con el suelo ($f_B$):} Como $B$ se mueve a la izquierda, el suelo rugoso raspa contra él hacia la derecha:
    \begin{equation}
        f_B = \mu \cdot N_B = 0.3 \cdot 60 \, \text{N} = 18 \, \text{N}
    \end{equation}
    \item \textbf{La fricción entre A y B ($f_{A \to B}$):} ¡Esta es la clave conceptual! Por la Tercera Ley de Newton (acción y reacción), si el bloque $B$ raspa al bloque $A$ con una fuerza de $6 \, \text{N}$ hacia la izquierda, el bloque $A$ reacciona raspando al bloque $B$ con exactamente $6 \, \text{N}$ hacia la derecha.
\end{itemize}

\subsubsection*{Equilibrio Final y Conclusión}
Para que el bloque $B$ se mueva con rapidez constante hacia la izquierda, la magnitud de la fuerza $\vec{F}$ debe ser capaz de igualar y vencer a las tres fuerzas que tiran hacia la derecha:
\begin{equation}
    F = T + f_B + f_{A \to B}
\end{equation}
\begin{equation}
    F = 6 \, \text{N} + 18 \, \text{N} + 6 \, \text{N}
\end{equation}
\begin{equation}
    F = 30 \, \text{N}
\end{equation}

La magnitud de la fuerza horizontal $\vec{F}$ necesaria para arrastrar el sistema a rapidez constante es de $30 \, \text{N}$.

\newpage

\subsection*{Enunciado}
Suponga que sostiene dos pelotas de tenis del mismo tamaño: una es una pelota regular hueca y la otra está llena de perdigones de hierro, lo cual la hace mucho más pesada. Si las suelta desde una altura pequeña (como la de sus brazos), observará que golpean el suelo casi al mismo tiempo. Explique por qué ocurre esto.

\subsection*{Solución}

\subsubsection*{Análisis de la Segunda Ley de Newton en Ausencia de Resistencia Significativa}
De acuerdo con los principios de la física conceptual, cuando un objeto se encuentra en caída libre ideal, la única fuerza que actúa sobre él es la interacción gravitacional. La segunda ley de Newton establece que la aceleración es directamente proporcional a la fuerza neta e inversamente proporcional a la masa ($a = \frac{F}{m}$). Dado que la fuerza peso de un objeto es $F = mg$, al sustituir esto en la ecuación obtenemos $a = \frac{mg}{m} = g$. Este paso es fundamental porque demuestra matemáticamente que la inercia (la masa en el denominador) anula exactamente el incremento de fuerza gravitacional (la masa en el numerador). Por lo tanto, la aceleración de caída es intrínsecamente independiente de la masa del objeto.

\subsubsection*{El Efecto de la Resistencia del Aire a Bajas Velocidades}
En el mundo real, la fuerza de resistencia del aire es una variable dependiente de la rapidez del objeto y de su perfil geométrico. Cuando las pelotas se sueltan desde una altura tan pequeña como la extensión de los brazos, el intervalo de tiempo de la caída es sumamente breve. Consecuentemente, las pelotas no tienen la oportunidad cinemática de adquirir una rapidez que genere una fuerza de fricción del aire con un valor que compita con su peso. Al ser esta fuerza de arrastre aerodinámico matemáticamente despreciable en comparación con el vector peso de ambas pelotas, la fuerza neta actuante es prácticamente equivalente a la fuerza gravitacional pura para ambos cuerpos.

\subsubsection*{Conclusión}
Debido a que a bajas velocidades y distancias cortas la resistencia del aire es mínima y se puede despreciar en el sistema, la fuerza neta sobre ambas pelotas es únicamente su peso. Esto provoca que ambas adquieran la misma aceleración constante (la aceleración de la gravedad $g$), sin importar que la pelota con perdigones posea mucha más masa. Por esta razón cinemática, ambas recorren la misma distancia en el mismo tiempo y golpean el suelo casi simultáneamente.

\newpage

\subsection*{Enunciado}
Suponga que sostiene dos pelotas de tenis del mismo tamaño: una es una pelota regular hueca y la otra está llena de perdigones de hierro, lo cual la hace mucho más pesada. Si las suelta desde la terraza de un edificio alto, notará que la pelota más pesada golpea el suelo primero. Utilice la segunda ley de Newton y el concepto de fuerza de resistencia ($\vec{R}$) para explicar por qué la pelota ligera experimenta una disminución de aceleración mucho más drástica que la pelota pesada, a pesar de encontrar la misma resistencia del aire a una rapidez dada.

\subsection*{Solución}

\subsubsection*{Definición de las Fuerzas Actuantes en Trayectos Largos}
Cuando las pelotas se sueltan desde una altura considerable, como la terraza de un edificio, el tiempo prolongado de caída les permite adquirir una rapidez significativa. A estas rapideces mayores, la resistencia del aire ($\vec{R}$) crece hasta convertirse en un factor que ya no se puede ignorar. Sobre cada pelota actúan ahora dos vectores de fuerza en direcciones opuestas: el vector peso hacia abajo ($W = mg$) y el vector de la resistencia del aire hacia arriba ($R$). Como ambas pelotas poseen idéntico tamaño y forma, sus áreas frontales son iguales. Por lo tanto, a cualquier rapidez instantánea dada $v$, ambas experimentan de forma garantizada exactamente la misma magnitud de fuerza de resistencia aerodinámica $R$.

\subsubsection*{Aplicación de la Segunda Ley de Newton en Fluidos}
Para determinar cómo se comporta cinemáticamente cada pelota ante la presencia del fluido (el aire), recurrimos nuevamente a la segunda ley de Newton ($a = \frac{F_{\text{neta}}}{m}$). La fuerza neta ahora se define como la diferencia vectorial entre el peso descendente y la resistencia ascendente. Esto genera la ecuación $a = \frac{mg - R}{m}$.
Para obtener una visualización pedagógica más profunda, se debe separar esta fracción algebraicamente en dos términos, obteniendo la siguiente expresión conceptual:
$a = g - \frac{R}{m}$

\subsubsection*{Análisis Matemático del Término de Retardo}
La ecuación $a = g - \frac{R}{m}$ revela el núcleo del fenómeno físico. La aceleración de caída ya no es la constante $g$, sino que sufre una penalización determinada por el cociente $\frac{R}{m}$. Para la pelota ligera (cuya masa $m$ es pequeña), dividir un valor $R$ entre un número pequeño da como resultado un valor relativamente grande. Esto implica que a la gravedad $g$ se le resta una cantidad sustancial, disminuyendo la aceleración drásticamente.
Por el contrario, para la pelota pesada llena de perdigones (cuya masa $m$ es muy grande), dividir ese mismo valor $R$ entre un número tan grande da como resultado una fracción minúscula. A la gravedad $g$ se le resta muy poco, permitiendo que la pelota pesada mantenga una aceleración muy cercana al valor de caída libre ideal a lo largo de su descenso.

\subsubsection*{Representación Gráfica de Diagramas de Cuerpo Libre}
Para evidenciar visualmente la desproporción entre la fuerza de resistencia compartida y los pesos individuales, se diseñará el siguiente modelado computacional que traza los vectores de fuerza para un instante donde ambas pelotas comparten la misma rapidez $v$.

\begin{lstlisting}[language=Python]
import matplotlib.pyplot as plt

fig, (ax1, ax2) = plt.subplots(1, 2, figsize=(9, 5))

ax1.set_xlim(-1, 1)
ax1.set_ylim(-12, 6)
ax1.axis('off')
ax1.set_title('Pelota Ligera (Hueca)', fontsize=14, fontweight='bold')
ax1.plot(0, 0, 'ro', markersize=25, markeredgecolor='black')
ax1.annotate('', xy=(0, 4), xytext=(0, 0), arrowprops=dict(facecolor='blue', width=3, headwidth=10))
ax1.text(0.1, 2, 'R', color='blue', fontsize=14, fontweight='bold')
ax1.annotate('', xy=(0, -5), xytext=(0, 0), arrowprops=dict(facecolor='red', width=3, headwidth=10))
ax1.text(0.1, -3, 'W = mg', color='red', fontsize=14, fontweight='bold')

ax2.set_xlim(-1, 1)
ax2.set_ylim(-12, 6)
ax2.axis('off')
ax2.set_title('Pelota Pesada (Perdigones)', fontsize=14, fontweight='bold')
ax2.plot(0, 0, 'go', markersize=25, markeredgecolor='black')
ax2.annotate('', xy=(0, 4), xytext=(0, 0), arrowprops=dict(facecolor='blue', width=3, headwidth=10))
ax2.text(0.1, 2, 'R', color='blue', fontsize=14, fontweight='bold')
ax2.annotate('', xy=(0, -11), xytext=(0, 0), arrowprops=dict(facecolor='darkred', width=6, headwidth=15))
ax2.text(0.1, -6, 'W = Mg', color='darkred', fontsize=14, fontweight='bold')

plt.tight_layout()
plt.show()
\end{lstlisting}

\subsubsection*{Conclusión}
A pesar de que ambas pelotas sufren la misma fuerza de arrastre $R$ por moverse a la misma rapidez, el impacto de esta resistencia en su estado de movimiento depende completamente de su inercia. Como el término penalizador $\frac{R}{m}$ es notablemente mayor para la pelota hueca debido a su poca masa, sufre una reducción de aceleración pronunciada, alcanzando su velocidad terminal mucho antes. En consecuencia, la pelota pesada logra mantener una aceleración más alta durante mayor tiempo, acumulando una rapidez media superior que le permite recorrer la distancia hasta el suelo en menor tiempo.

\newpage

\subsection*{Enunciado}
Un paracaidista de 80 kg salta desde un avión y alcanza una rapidez terminal de 50 m/s con los brazos y piernas extendidos. Determine el valor de la fuerza de arrastre ($R$) que actúa sobre el paracaidista cuando ha alcanzado su rapidez terminal. Explique su respuesta basándose en la aceleración del objeto en ese punto.

\subsection*{Solución}

\subsubsection*{Resultado Final}
El valor de la fuerza de arrastre $R$ es de $784\text{ N}$.

\subsubsection*{Análisis del Sistema}
La lógica inicial y los parámetros de masa y rapidez terminal son validados. De acuerdo con los principios de la mecánica de fluidos y la segunda ley de Newton, la condición de rapidez terminal establece un estado de equilibrio dinámico donde la aceleración del sistema es nula.

\subsubsection*{Transformación Matemática}
Se procesa el balance de fuerzas asumiendo gravedad estándar $g = 9.8\text{ m/s}^2$.
$R = mg = 80 \cdot 9.8 = 784\text{ N}$

\subsubsection*{Conclusión}
La fuerza de arrastre $R$ que actúa sobre el paracaidista en su rapidez terminal es exactamente $784\text{ N}$.

\subsection*{Enunciado}
Un paracaidista de 80 kg salta desde un avión y alcanza una rapidez terminal de 50 m/s con los brazos y piernas extendidos. Determine la magnitud de la aceleración del paracaidista en el instante en que su rapidez es de 30 m/s. Para ello, asuma que la fuerza de resistencia sigue el modelo $R = Dv^2$, donde $D$ es una constante de arrastre que puede despejar de los datos de rapidez terminal.

\subsection*{Solución}

\subsubsection*{Resultado Final}
La magnitud de la aceleración es $a = 6.272\text{ m/s}^2$.

\subsubsection*{Análisis del Sistema}
La lógica inicial y los parámetros cinemáticos son validados. La constante aerodinámica $D$ se extrae de la condición límite de velocidad terminal previamente establecida y se aplica al modelo de resistencia cuadrática de forma directa.

\subsubsection*{Transformación Matemática}
Se transfiere la constante $D$ a la ecuación diferencial de movimiento en el instante $v = 30\text{ m/s}$.
$a = g - \frac{D(30)^2}{m} = 6.272\text{ m/s}^2$

\subsubsection*{Conclusión}
La aceleración del paracaidista en el instante requerido es $6.272\text{ m/s}^2$.

\newpage

\subsection*{Enunciado}
¿Cuál es el origen físico de la fuerza de fricción entre dos superficies sólidas en contacto?

a) La presión atmosférica que empuja las superficies entre sí.

b) Únicamente el peso del objeto que se desliza.

c) Las irregularidades microscópicas y los enlaces químicos entre átomos en los puntos de contacto.

d) La repulsión magnética entre los electrones de las superficies.

e) El movimiento térmico de las moléculas en la interfaz de contacto.

\subsection*{Solución}

\subsubsection*{Análisis de las Superficies a Nivel Microscópico}
De acuerdo con la aproximación de la física conceptual para la mecánica de contacto, es fundamental entender que ninguna superficie sólida macroscópica es perfectamente lisa. Incluso en los materiales que parecen estar pulidos a la perfección a simple vista, la observación microscópica revela un relieve sumamente complejo e irregular, compuesto por montañas y valles a nivel atómico. Esta naturaleza rugosa subyacente es el pilar para comprender cómo interactúan los materiales.

\subsubsection*{Interacción en los Puntos de Contacto Real}
Cuando se coloca un objeto sólido sobre otro, el área de apoyo real es en realidad una fracción minúscula del área que aparentan tener en contacto. Las dos superficies únicamente se tocan en los picos más altos de sus respectivas irregularidades microscópicas. En estos minúsculos puntos de contacto íntimo, la distancia que separa a los átomos del objeto superior y los del inferior es tan reducida que las fuerzas electromagnéticas intermoleculares entran en acción. Esto provoca que los electrones interactúen, formando enlaces químicos temporales (fuerzas de adhesión o soldaduras frías) entre las dos superficies.

\subsubsection*{Justificación del Origen de la Resistencia al Deslizamiento}
La fuerza de fricción estática o cinética es la manifestación macroscópica del esfuerzo necesario para vencer dos barreras físicas al intentar iniciar o mantener un deslizamiento. Primero, se requiere energía para fracturar o romper de manera continua esos enlaces químicos formados en los puntos de apoyo. Segundo, se necesita fuerza para deformar, arrancar o hacer saltar las irregularidades (los picos) que se han entrelazado mecánicamente entre sí. 
Otras variables mencionadas en las alternativas tienen roles diferentes: el peso del objeto (relacionado con la fuerza normal) simplemente aumenta la presión sobre estos picos incrementando el área real de contacto y la cantidad de enlaces, pero no es el origen del fenómeno en sí; la presión atmosférica o el movimiento térmico no definen mecánicamente la resistencia al corte en la interfaz.

\subsubsection*{Conclusión}
La opción correcta es la c. La fricción surge fundamentalmente debido al entrelazamiento físico de las irregularidades microscópicas presentes en los materiales y a la constante formación y ruptura de enlaces químicos entre los átomos en los verdaderos puntos de contacto de ambas superficies.

\newpage

\subsection*{Enunciado}
Si se aplica una fuerza horizontal externa $\vec{F}$ a un objeto en reposo sobre una superficie horizontal rugosa y este no se mueve, la magnitud de la fuerza de fricción estática es:

a) $fr_e = \mu_e N$

b) $fr_e = F$

c) $fr_e < F$

d) $fr_e > F$

e) $fr_e = 0$

\subsection*{Solución}

\subsubsection*{Análisis de la Primera Ley de Newton y el Equilibrio Estático}
El punto de partida para resolver esta situación es la observación cinemática de que el objeto permanece en reposo. Según la primera ley de Newton (ley de la inercia) y la segunda ley ($\Sigma F = ma$), si un objeto no presenta aceleración ($a = 0$), la suma vectorial de todas las fuerzas que actúan sobre él debe ser estrictamente cero. Esto significa que el objeto se encuentra en un estado de equilibrio mecánico. En el eje horizontal, la fuerza neta debe anularse por completo.

\subsubsection*{La Naturaleza Reactiva de la Fricción Estática}
Cuando se aplica la fuerza horizontal externa $F$, el objeto tiende a moverse. La superficie rugosa responde ejerciendo una fuerza paralela a la superficie y en sentido opuesto al movimiento inminente, conocida como fuerza de fricción estática ($fr_e$). 
Es un error común confundir la fuerza de fricción estática real con su valor máximo posible. La fórmula $fr_{e,\text{máx}} = \mu_e N$ (opción a) describe únicamente el límite de ruptura, es decir, la fuerza máxima que la superficie puede soportar antes de que el objeto comience a deslizar. Sin embargo, mientras el objeto no se mueva, la fricción estática actúa como una fuerza "inteligente" o reactiva: ajusta su magnitud exactamente para igualar y cancelar la fuerza aplicada $F$. Si la fuerza aplicada aumenta sin superar el límite, la fricción estática aumenta en la misma proporción para mantener el equilibrio. Por lo tanto, para garantizar que la suma de fuerzas en el eje horizontal sea cero ($\Sigma F_x = F - fr_e = 0$), es matemáticamente imperativo que $fr_e = F$.

\subsubsection*{Representación Gráfica del Diagrama de Cuerpo Libre}
Para ilustrar este equilibrio de fuerzas, el siguiente modelado computacional muestra cómo los vectores horizontales y verticales se cancelan entre sí, garantizando el reposo del sistema.

\begin{lstlisting}[language=Python]
import matplotlib.pyplot as plt

fig, ax = plt.subplots(figsize=(7, 5))
ax.set_xlim(-5, 5)
ax.set_ylim(-4, 5)
ax.axis('off')

ax.plot([-4, 4], [0, 0], color='black', linewidth=2)

rect = plt.Rectangle((-1.5, 0), 3, 2, facecolor='lightgray', edgecolor='black', linewidth=1.5)
ax.add_patch(rect)

ax.plot(0, 1, 'ko', markersize=5)

ax.annotate('', xy=(3.5, 1), xytext=(1.5, 1), arrowprops=dict(facecolor='blue', width=3, headwidth=10))
ax.text(2.5, 1.4, 'F (Aplicada)', color='blue', fontsize=12, fontweight='bold', ha='center')

ax.annotate('', xy=(-3.5, 0), xytext=(-1.5, 0), arrowprops=dict(facecolor='red', width=3, headwidth=10))
ax.text(-2.5, -0.6, 'fre (Friccion)', color='red', fontsize=12, fontweight='bold', ha='center')

ax.annotate('', xy=(0, 4), xytext=(0, 2), arrowprops=dict(facecolor='green', width=3, headwidth=10))
ax.text(0.2, 3, 'N', color='green', fontsize=12, fontweight='bold')

ax.annotate('', xy=(0, -3), xytext=(0, 0), arrowprops=dict(facecolor='purple', width=3, headwidth=10))
ax.text(0.2, -2, 'W = mg', color='purple', fontsize=12, fontweight='bold')

plt.title('Diagrama de Cuerpo Libre: Equilibrio Estatico', fontsize=14, fontweight='bold')
plt.tight_layout()
plt.show()
\end{lstlisting}

\subsubsection*{Conclusión}
La opción correcta es la b. Mientras el objeto permanezca en reposo bajo la acción de una fuerza externa, el sistema se encuentra en equilibrio traslacional. Para satisfacer la primera ley de Newton, la fuerza de fricción estática generada por la superficie debe igualar exactamente la magnitud de la fuerza aplicada ($fr_e = F$) para evitar que exista una fuerza neta y, por ende, una aceleración.

\newpage

\subsection*{Enunciado}
Si se empuja horizontalmente una caja que descansa sobre un piso nivelado con una fuerza de 70 N y esta no se mueve, ¿cuál es la magnitud de la fuerza de fricción?

a) Es un poco mayor a 70 N.

b) Es exactamente igual a 70 N.

c) Es igual al coeficiente de fricción estática multiplicado por el peso de la caja.

d) Cero, porque el objeto no está en movimiento.

e) Depende de la rapidez con la que se aplicó la fuerza.

\subsection*{Solución}

\subsubsection*{Análisis del Sistema}
La lógica inicial y el parámetro de fuerza externa aplicada son validados. De acuerdo con los principios de la primera ley de Newton, la ausencia de movimiento en la caja confirma un estado de equilibrio traslacional donde la sumatoria de todas las fuerzas actuantes en el eje horizontal es estrictamente nula.

\subsubsection*{Transformación Matemática}
Se procesa el balance de fuerzas horizontales para despejar la magnitud reactiva.
$\Sigma F_x = F_{\text{aplicada}} - f_{r} = 0 \Rightarrow f_{r} = 70 \text{ N}$

\subsubsection*{Conclusión}
La opción correcta es la b. La fuerza de fricción estática ajusta su magnitud para ser exactamente 70 N en sentido opuesto a la fuerza aplicada, garantizando así que la fuerza neta sea cero y el objeto se mantenga en reposo.

\newpage

\subsection*{Enunciado}
A diferencia de la fricción entre superficies sólidas (seca), la fuerza de resistencia en un fluido (líquido o gas):

a) No depende del área transversal del objeto.

b) Actúa en la misma dirección que el vector velocidad.

c) Siempre es constante independientemente de la rapidez.

d) Depende fundamentalmente de la rapidez del objeto a través del medio.

e) Solo existe cuando el objeto alcanza la rapidez terminal.

\subsection*{Solución}

\subsubsection*{Diferenciación entre Fricción Seca y Resistencia en Fluidos}
Para comprender esta distinción, es crucial basarse en los modelos físicos descritos en la obra de Hewitt. La fricción cinética entre dos superficies sólidas ocurre debido a las interacciones de adhesión y entrelazamiento de las asperezas microscópicas. Una característica notable de este tipo de fricción es que, en primera aproximación, es casi independiente de la rapidez con la que se deslicen las superficies una sobre otra.

Por el contrario, la resistencia en un fluido (denominada arrastre aerodinámico o hidrodinámico) no se basa en el contacto mecánico de asperezas, sino en la necesidad del objeto de desplazar y empujar a las moléculas del fluido para abrirse paso. A mayor rapidez, el objeto debe desplazar más masa de fluido por unidad de tiempo y con mayor intensidad.

\subsubsection*{Dependencia Física con la Rapidez}
La fuerza de resistencia ($R$) en un fluido está intrínsecamente ligada a la rapidez ($v$). Según modelos estándar de mecánica, esta fuerza se describe frecuentemente como $R \propto v$ para regímenes de flujo laminar (ley de Stokes) o $R \propto v^2$ para regímenes de flujo turbulento (objetos moviéndose a altas rapideces en aire, como paracaidistas). Esta dependencia cuadrática o lineal con la rapidez es la razón fundamental por la cual la fuerza de arrastre aumenta drásticamente cuando la velocidad aumenta, a diferencia de la fricción sólida que permanece relativamente estable.

\subsubsection*{Análisis de las Otras Alternativas}
Las demás opciones son incorrectas por los siguientes motivos:
\begin{itemize}
    \item La resistencia sí depende fuertemente del área transversal (a).
    \item La resistencia siempre actúa en dirección opuesta al vector velocidad, no en la misma (b).
    \item Como se explicó, la resistencia no es constante; varía exponencialmente con la rapidez (c).
    \item La resistencia existe desde el momento en que el objeto comienza a moverse en el fluido, no solo en la rapidez terminal (e).
\end{itemize}


\subsubsection*{Conclusión}
La opción correcta es la d. A diferencia de la fricción cinética entre sólidos, la fuerza de resistencia que un fluido ejerce sobre un objeto depende fundamentalmente de la rapidez de dicho objeto a través del medio.

\newpage

\subsection*{Enunciado}
¿Qué sucede con un paracaidista inmediatamente después de abrir su paracaídas tras haber alcanzado su primera rapidez terminal?

a) Empieza a subir físicamente hacia el cielo.

b) La resistencia del aire aumenta drásticamente, volviéndose mayor que el peso y causando una aceleración hacia arriba.

c) La fuerza neta sigue siendo nula, por lo que su velocidad no cambia.

d) La gravedad aumenta para compensar el área extra del paracaídas.

e) Alcanza una nueva rapidez terminal que es mayor que la anterior.

\subsection*{Solución}

\subsubsection*{Análisis del Equilibrio Previo a la Apertura}
De acuerdo con los principios de la física conceptual, cuando el paracaidista alcanza su primera rapidez terminal cayendo en caída libre con su cuerpo extendido, se encuentra en un estado de equilibrio dinámico. En este punto, la fuerza de resistencia del aire empujando hacia arriba es exactamente igual en magnitud al peso del paracaidista tirando hacia abajo ($R = mg$). Como la fuerza neta es cero, la aceleración es nula y la velocidad de caída se mantiene constante.

\subsubsection*{El Efecto del Incremento Drástico del Área Frontal}
Inmediatamente al desplegar el paracaídas, el área transversal expuesta al fluido (el aire) aumenta de forma masiva. Dado que la fuerza de resistencia aerodinámica depende directamente de esta área frontal y del coeficiente de arrastre, el valor de la resistencia experimenta un incremento drástico e instantáneo. Como el paracaidista aún lleva la alta velocidad de su primera rapidez terminal, esta nueva fuerza de resistencia aerodinámica se vuelve significativamente mayor que la fuerza de gravedad constante (su peso). Matemáticamente, la fuerza neta sobre el sistema es ahora ascendente ($F_{\text{neta}} = R - mg > 0$).

\subsubsection*{Interpretación Cinemática de la Aceleración Hacia Arriba}
La segunda ley de Newton establece que una fuerza neta ascendente produce obligatoriamente una aceleración en la misma dirección (hacia arriba). Es en este punto donde se debe desmentir el error conceptual visualizado comúnmente desde el marco de referencia de una cámara en caída libre: el paracaidista no "sube" físicamente en relación al suelo. Dado que el vector velocidad del paracaidista sigue apuntando hacia abajo, una aceleración que apunta en sentido contrario (hacia arriba) actúa puramente como una fuerza de frenado. Esta aceleración ascendente reduce de forma abrupta la magnitud de la rapidez de caída, forzando una desaceleración hasta que se alcance una nueva y menor rapidez terminal.

\subsubsection*{Representación Gráfica del Desequilibrio Transitorio}
Para visualizar esta interacción vectorial en el instante exacto del despliegue del paracaídas, se plantea el siguiente modelado computacional que contrasta las magnitudes relativas de ambas fuerzas.

\begin{lstlisting}[language=Python]
import matplotlib.pyplot as plt

fig, ax = plt.subplots(figsize=(6, 8))
ax.set_xlim(-3, 3)
ax.set_ylim(-6, 12)
ax.axis('off')

ax.plot(0, 0, 'ko', markersize=15)

ax.annotate('', xy=(0, 10), xytext=(0, 0), arrowprops=dict(facecolor='blue', width=6, headwidth=15))
ax.text(0.5, 5, 'R (Nueva Resistencia)', color='blue', fontsize=14, fontweight='bold')

ax.annotate('', xy=(0, -4), xytext=(0, 0), arrowprops=dict(facecolor='red', width=3, headwidth=10))
ax.text(0.5, -2, 'W = mg', color='red', fontsize=14, fontweight='bold')

plt.title('Fuerzas Inmediatamente Despues de Abrir el Paracaidas', fontsize=12, fontweight='bold')
plt.tight_layout()
plt.show()
\end{lstlisting}

\subsubsection*{Conclusión}
La opción correcta es la b. La apertura del paracaídas provoca que la resistencia del aire supere abrumadoramente al peso. Esta fuerza neta ascendente genera una aceleración dirigida hacia arriba que, al ser opuesta al vector de velocidad descendente, frena de manera contundente la caída del paracaidista sin causar en ningún momento que su cuerpo se desplace físicamente hacia el cielo.

\newpage

\subsection*{Enunciado}
Un disco de hockey sobre un lago congelado recibe una rapidez inicial de 20 m/s. El disco se desliza una distancia de 115 m antes de llegar al reposo. Determine el coeficiente de fricción cinética ($\mu_c$) entre el disco y el hielo. Asuma una aceleración de la gravedad $g = 10 \text{ m/s}^2$.

\subsection*{Solución}

\subsubsection*{Análisis Cinemático del Deslizamiento}
Para relacionar las fuerzas con el movimiento del disco, primero debemos entender su comportamiento cinemático. El disco inicia con una rapidez de 20 m/s y se detiene completamente tras recorrer 115 m, lo que implica una aceleración negativa constante que frena el objeto. Utilizando las ecuaciones del movimiento uniformemente variado, específicamente la que relaciona las velocidades, la aceleración y el desplazamiento horizontal ($v_{xf}^2 = v_{x0}^2 + 2a_x\Delta x$), podemos despejar la aceleración del sistema.
$0 = (20)^2 + 2a_x(115)$
$a_x = \frac{-400}{230} = -1.74 \text{ m/s}^2$

\subsubsection*{Identificación de Fuerzas y Diagrama de Cuerpo Libre}
Según la física conceptual, un objeto en movimiento sobre una superficie interactúa mecánicamente con su entorno. Sobre el disco actúan tres fuerzas fundamentales: la interacción gravitacional hacia abajo (el peso, $P = mg$), la fuerza de soporte o contacto de la superficie de hielo hacia arriba (la fuerza normal, $N$), y la fuerza resistente que se opone al deslizamiento (la fricción cinética, $fr_c$). Para visualizar esto con claridad, el siguiente modelado computacional grafica el diagrama de cuerpo libre del disco.

\begin{lstlisting}[language=Python]
import matplotlib.pyplot as plt

fig, ax = plt.subplots(figsize=(5, 5))
ax.set_xlim(-4, 4)
ax.set_ylim(-4, 4)
ax.axis('off')

rect = plt.Rectangle((-1, -1), 2, 2, facecolor='whitesmoke', edgecolor='black', linewidth=1.5)
ax.add_patch(rect)

ax.annotate('', xy=(0, 3), xytext=(0, 1), arrowprops=dict(facecolor='black', width=2, headwidth=8))
ax.text(0.2, 2.5, 'N', fontsize=14, fontweight='bold')

ax.annotate('', xy=(0, -3), xytext=(0, -1), arrowprops=dict(facecolor='black', width=2, headwidth=8))
ax.text(0.2, -2.5, 'P = mg', fontsize=14, fontweight='bold')

ax.annotate('', xy=(-3, 0), xytext=(-1, 0), arrowprops=dict(facecolor='black', width=2, headwidth=8))
ax.text(-2.5, -0.5, 'frc', fontsize=14, fontweight='bold')

plt.title('DCL: Disco', fontsize=14, fontweight='bold')
plt.tight_layout()
plt.show()
\end{lstlisting}

\subsubsection*{Aplicación de la Segunda Ley de Newton}
En el eje vertical (y), el disco no presenta movimiento hacia arriba ni hacia abajo, por lo que se encuentra en perfecto equilibrio traslacional. La sumatoria de fuerzas es cero ($\Sigma F_y = 0$), lo que nos permite establecer de forma concluyente que la magnitud de la fuerza normal es exactamente igual al peso del disco ($N = mg$).
En el eje horizontal (x), la única fuerza neta que actúa sobre el disco es la fricción cinética, apuntando en dirección opuesta al movimiento. Al aplicar la segunda ley de Newton en este eje ($\Sigma F_x = ma_x$), establecemos que la fuerza neta es $-fr_c = ma_x$.

\subsubsection*{Relación de la Fricción y el Coeficiente Cinético}
Físicamente, la magnitud de la fuerza de fricción cinética es directamente proporcional a la fuerza de soporte (la normal) que presiona las superficies entre sí, definida por la ecuación empírica $fr_c = \mu_c N$. Al sustituir la equivalencia de la normal obtenida anteriormente, llegamos a $fr_c = \mu_c mg$. 
Si introducimos esta nueva expresión en la ecuación de la segunda ley de Newton para el eje x, obtenemos:
$-\mu_c mg = ma_x$
La masa ($m$) se cancela en ambos lados de la igualdad, lo cual demuestra algebraicamente que la tasa de frenado depende de la rugosidad de las superficies, pero es totalmente independiente de la masa del disco. Despejando el coeficiente de fricción cinética ($\mu_c$) y utilizando el valor de aceleración cinemática previamente calculado:
$\mu_c = \frac{-a_x}{g}$
$\mu_c = \frac{-(-1.74)}{10}$

\subsubsection*{Conclusión}
El coeficiente de fricción cinética $\mu_c$ entre el disco y el hielo es de 0.174.

\newpage

\subsection*{Enunciado}
Se coloca un bloque sobre una rampa. Al inclinar la rampa lentamente, se detecta que el bloque comienza a deslizarse justo cuando el ángulo alcanza un valor $\theta_c$. Si el coeficiente de fricción estática es $\mu_e = 0.350$, encuentre el ángulo crítico $\theta_c$.

\subsection*{Solución}

\subsubsection*{Análisis del Equilibrio y Sistema de Referencia}
Desde la perspectiva de la física conceptual, el instante preciso justo antes de que el bloque comience a deslizarse se denomina estado de movimiento inminente. En este punto crítico, el bloque aún se encuentra en reposo relativo a la rampa, por lo que cumple estrictamente con la primera condición de equilibrio traslacional ($\Sigma\vec{F} = \vec{0}$). Para simplificar el análisis matemático, es pedagógicamente útil rotar nuestro sistema de coordenadas espaciales: alineamos el eje $x$ paralelo a la superficie de la rampa (dirección del posible deslizamiento) y el eje $y$ perpendicular a ella.

\subsubsection*{Diagrama de Cuerpo Libre}
Sobre el bloque actúan tres vectores de fuerza fundamentales: la fuerza de gravedad tirando verticalmente hacia abajo (peso $\vec{P}$), la superficie empujando hacia afuera (fuerza normal $\vec{N}$), y el agarre microscópico de las superficies oponiéndose a la caída (fuerza de fricción estática máxima $\vec{fr}_{em\acute{a}x}$).

\begin{lstlisting}[language=Python]
import matplotlib.pyplot as plt
import numpy as np
from matplotlib.patches import Polygon

fig, ax = plt.subplots(figsize=(6, 6))
ax.set_xlim(-4, 4)
ax.set_ylim(-4, 4)
ax.axis('off')

theta = np.radians(20)
c, s = np.cos(theta), np.sin(theta)

w, h = 2, 1.5
pts = np.array([[-w/2, -h/2], [w/2, -h/2], [w/2, h/2], [-w/2, h/2]])
rot_pts = np.zeros_like(pts)
for i in range(4):
    x, y = pts[i]
    rot_pts[i, 0] = x*c + y*s
    rot_pts[i, 1] = -x*s + y*c

block = Polygon(rot_pts, closed=True, facecolor='whitesmoke', edgecolor='black', linewidth=1.5)
ax.add_patch(block)

nx, ny = -1.5*s, 1.5*c
ax.annotate('', xy=(nx, ny), xytext=(0, 0), arrowprops=dict(facecolor='black', width=2, headwidth=8))
ax.text(nx*1.1, ny*1.1, 'N', fontsize=14, fontweight='bold', ha='center')

ax.annotate('', xy=(0, -2.5), xytext=(0, 0), arrowprops=dict(facecolor='black', width=2, headwidth=8))
ax.text(0.2, -2.8, 'P = mg', fontsize=14, fontweight='bold')

fx, fy = -1.8*c, 1.8*s
ax.annotate('', xy=(fx, fy), xytext=(0, 0), arrowprops=dict(facecolor='black', width=2, headwidth=8))
ax.text(fx*1.2, fy*1.2, 'fre_max', fontsize=14, fontweight='bold', ha='right')

ax.annotate('', xy=(3*c, -3*s), xytext=(2*c, -2*s), arrowprops=dict(arrowstyle="->", color='gray'))
ax.text(3.2*c, -3.2*s, 'x', color='gray', fontsize=12)
ax.annotate('', xy=(2*c + 1*s, -2*s + 1*c), xytext=(2*c, -2*s), arrowprops=dict(arrowstyle="->", color='gray'))
ax.text((2*c + 1.2*s), (-2*s + 1.2*c), 'y', color='gray', fontsize=12)

plt.title('DCL: Movimiento Inminente', fontsize=14, fontweight='bold')
plt.tight_layout()
plt.show()
\end{lstlisting}

\subsubsection*{Descomposición de Fuerzas y Leyes de Newton}
Dado que el peso ($mg$) apunta directamente hacia el centro de la Tierra y no está alineado con nuestros ejes inclinados, debemos descomponerlo trigonométricamente. La componente que empuja el bloque contra la rampa es $mg\cos(\theta_c)$, y la componente paralela que intenta deslizar el bloque hacia abajo es $mg\text{sen}(\theta_c)$.

Aplicando el equilibrio en el eje $y$ (perpendicular a la rampa):
$$\Sigma F_y = N - mg\cos(\theta_c) = 0$$
$$N = mg\cos(\theta_c)$$

Aplicando el equilibrio en el eje $x$ (paralelo a la rampa):
$$\Sigma F_x = mg\text{sen}(\theta_c) - fr_{em\acute{a}x} = 0$$

\subsubsection*{Condición de Movimiento Inminente y Cancelación de la Masa}
Sabemos que la fuerza de fricción estática alcanza su límite máximo absoluto justo antes del deslizamiento, definido empíricamente como el producto del coeficiente de fricción y la fuerza normal presionante.
$$fr_{em\acute{a}x} = \mu_e N$$

Sustituyendo el valor de la normal obtenido del eje $y$:
$$fr_{em\acute{a}x} = \mu_e mg\cos(\theta_c)$$

Ahora, insertamos esta expresión en la ecuación de equilibrio del eje $x$:
$$mg\text{sen}(\theta_c) = \mu_e mg\cos(\theta_c)$$

En este paso ocurre una simplificación conceptual fascinante: la masa ($m$) y la aceleración de la gravedad ($g$) se encuentran en ambos lados de la igualdad y pueden cancelarse. Esto demuestra matemáticamente que el ángulo al que un objeto comienza a resbalar es totalmente independiente de su peso. Una caja fuerte muy pesada y un bloque ligero de madera comenzarán a deslizarse exactamente al mismo ángulo de inclinación si sus coeficientes de fricción son iguales.
Despejando el coeficiente obtenemos:
$$\frac{\text{sen}(\theta_c)}{\cos(\theta_c)} = \mu_e$$
$$\tan(\theta_c) = \mu_e$$

\subsubsection*{Cálculo del Ángulo Crítico}
Sustituyendo el valor numérico del coeficiente de fricción estática ($\mu_e = 0.350$) proporcionado en el problema:
$$\theta_c = \tan^{-1}(0.350)$$
$$\theta_c = 19.3^\circ$$

\subsubsection*{Conclusión}
El ángulo crítico $\theta_c$ al cual el bloque comienza a deslizarse por la rampa es de $19.3^\circ$.

\newpage

\subsection*{Enunciado}
Un bloque de 3 kg parte del reposo y se desliza 2 m hacia abajo por un plano inclinado 30$^\circ$ respecto a la horizontal, en un tiempo de 1.5 s. Determine la magnitud de la aceleración del bloque.

\subsection*{Solución}

\subsubsection*{Análisis Cinemático del Movimiento}
Para abordar el problema desde la física conceptual, primero debemos analizar el estado de movimiento del bloque. Al indicar que "parte del reposo", se establece que la velocidad inicial es nula ($v_{x0} = 0$). Dado que el bloque se desliza sobre una superficie con una inclinación constante y asumiendo que las fuerzas resistivas no varían, la fuerza neta es constante. Según la segunda ley de Newton, una fuerza neta constante produce una aceleración constante. Este escenario nos permite utilizar las ecuaciones del movimiento rectilíneo uniformemente acelerado para describir su trayectoria espacial y temporal.

\subsubsection*{Cálculo de la Aceleración}
Utilizaremos la ecuación cinemática que relaciona el desplazamiento ($\Delta x$), la velocidad inicial ($v_{x0}$), el tiempo ($\Delta t$) y la aceleración ($a_x$):
$$\Delta x = v_{x0}\Delta t + \frac{1}{2}a_x\Delta t^2$$
Sustituyendo los valores numéricos proporcionados por el sistema ($v_{x0} = 0$, $\Delta x = 2 \text{ m}$, $\Delta t = 1.5 \text{ s}$):
$$2 = 0 + \frac{1}{2}a_x(1.5)^2$$
Despejando algebraicamente la variable de aceleración $a_x$:
$$a_x = \frac{2(2)}{(1.5)^2}$$
$$a_x = \frac{4}{2.25} = 1.78 \text{ m/s}^2$$

\subsubsection*{Conclusión}
La magnitud de la aceleración constante que experimenta el bloque al descender por el plano inclinado es de $1.78 \text{ m/s}^2$.

\newpage

\subsection*{Enunciado}
Un bloque de 3 kg parte del reposo y se desliza 2 m hacia abajo por un plano inclinado 30$^\circ$ respecto a la horizontal, en un tiempo de 1.5 s. Determine la magnitud de la fuerza de fricción cinética que actúa sobre el bloque. Asuma una aceleración de la gravedad $g = 10 \text{ m/s}^2$.

\subsection*{Solución}

\subsubsection*{Identificación de Fuerzas y Sistema de Referencia}
Para determinar las fuerzas interactuantes, es indispensable definir un sistema de coordenadas conveniente. Alineamos el eje $x$ paralelamente al plano inclinado (hacia abajo, en dirección del movimiento) y el eje $y$ perpendicular a este. Las fuerzas que actúan sobre el bloque son su peso ($\vec{P}$), que apunta verticalmente hacia el centro de la Tierra; la fuerza normal ($\vec{N}$), perpendicular a la superficie de contacto; y la fuerza de fricción cinética ($\vec{fr}_c$), que actúa paralela a la superficie y en sentido opuesto al deslizamiento.

\begin{lstlisting}[language=Python]
import matplotlib.pyplot as plt
import numpy as np
from matplotlib.patches import Polygon

fig, ax = plt.subplots(figsize=(6, 6))
ax.set_xlim(-4, 4)
ax.set_ylim(-4, 4)
ax.axis('off')

theta = np.radians(30)
c, s = np.cos(theta), np.sin(theta)

w, h = 1.5, 1.0
pts = np.array([[-w/2, -h/2], [w/2, -h/2], [w/2, h/2], [-w/2, h/2]])
rot_pts = np.zeros_like(pts)
for i in range(4):
    x, y = pts[i]
    rot_pts[i, 0] = x*c + y*s
    rot_pts[i, 1] = -x*s + y*c

block = Polygon(rot_pts, closed=True, facecolor='whitesmoke', edgecolor='black', linewidth=1.5)
ax.add_patch(block)

nx, ny = -1.2*s, 1.2*c
ax.annotate('', xy=(nx, ny), xytext=(0, 0), arrowprops=dict(facecolor='black', width=2, headwidth=8))
ax.text(nx*1.2, ny*1.2, 'N', fontsize=14, fontweight='bold', ha='center')

ax.annotate('', xy=(0, -2.5), xytext=(0, 0), arrowprops=dict(facecolor='black', width=2, headwidth=8))
ax.text(0.2, -2.8, 'P = mg', fontsize=14, fontweight='bold')

fx, fy = -1.8*c, 1.8*s
ax.annotate('', xy=(fx, fy), xytext=(0, 0), arrowprops=dict(facecolor='black', width=2, headwidth=8))
ax.text(fx*1.2, fy*1.2, 'frc', fontsize=14, fontweight='bold', ha='right')

ax.annotate('', xy=(3*c, -3*s), xytext=(2*c, -2*s), arrowprops=dict(arrowstyle="->", color='gray'))
ax.text(3.2*c, -3.2*s, 'x', color='gray', fontsize=12)
ax.annotate('', xy=(2*c + 1*s, -2*s + 1*c), xytext=(2*c, -2*s), arrowprops=dict(arrowstyle="->", color='gray'))
ax.text((2*c + 1.2*s), (-2*s + 1.2*c), 'y', color='gray', fontsize=12)

plt.title('DCL: Bloque en Plano Inclinado', fontsize=14, fontweight='bold')
plt.tight_layout()
plt.show()
\end{lstlisting}

\subsubsection*{Aplicación de la Segunda Ley de Newton}
El peso del bloque debe descomponerse. La componente responsable de acelerar el bloque hacia abajo por la rampa es $mg\text{sen}(\theta)$. La fuerza neta en el eje $x$ es la diferencia entre esta componente motriz y la fuerza de fricción resistente que se opone al movimiento. Aplicando la segunda ley de Newton ($\Sigma F_x = ma_x$), obtenemos la ecuación rectora del sistema dinámico:
$$\Sigma F_x = mg\text{sen}(\theta) - fr_c = ma_x$$

\subsubsection*{Cálculo de la Magnitud de la Fricción Cinética}
Para encontrar el valor de la fricción cinética, se manipula la ecuación anterior para despejar $fr_c$. La aceleración $a_x = 1.78 \text{ m/s}^2$ ya es un valor conocido extraído del análisis cinemático previo.
$$fr_c = mg\text{sen}(\theta) - ma_x$$
Factorizando la masa ($m$) para simplificar la expresión:
$$fr_c = m(g\text{sen}(\theta) - a_x)$$
Sustituyendo los valores de masa ($m = 3 \text{ kg}$), gravedad ($g = 10 \text{ m/s}^2$) y el ángulo ($\theta = 30^\circ$):
$$fr_c = 3(10\text{sen}(30^\circ) - 1.78)$$
Dado que $\text{sen}(30^\circ) = 0.5$:
$$fr_c = 3(5 - 1.78)$$
$$fr_c = 3(3.22) = 9.66 \text{ N}$$

\subsubsection*{Conclusión}
La magnitud de la fuerza de fricción cinética que se opone al deslizamiento del bloque es de $9.66 \text{ N}$.

\newpage

\subsection*{Enunciado}
Una pequeña esfera de 2 g se libera desde el reposo en aceite. El medio ejerce una fuerza resistiva proporcional a la rapidez ($R = bv$). Si la rapidez terminal es $v_T = 5 \text{ cm/s}$, determine la constante $b$. Asuma la aceleración de la gravedad como $g = 10 \text{ m/s}^2$.

\subsection*{Solución}

\subsubsection*{El Concepto de Rapidez Terminal en Fluidos}
Desde el enfoque de la física conceptual, cuando un objeto cae a través de un fluido viscoso como el aceite, no experimenta una caída libre ideal. Inicialmente, la fuerza de gravedad acelera la esfera hacia abajo. Sin embargo, a medida que su rapidez ($v$) aumenta, la fuerza resistiva del fluido ($R = bv$) que empuja hacia arriba crece proporcionalmente. El concepto clave aquí es la "rapidez terminal": es el instante físico en el cual la fuerza resistiva de arrastre crece lo suficiente como para igualar exactamente el peso de la esfera. En este punto, las fuerzas se cancelan, la fuerza neta se vuelve cero y, por la primera ley de Newton, la aceleración desaparece, permitiendo que la esfera continúe cayendo a una velocidad constante.

\subsubsection*{Diagrama de Cuerpo Libre y Equilibrio}
Para modelar esta condición de equilibrio dinámico, trazamos las fuerzas concurrentes sobre la esfera.

\begin{lstlisting}[language=Python]
import matplotlib.pyplot as plt
import matplotlib.patches as patches

fig, ax = plt.subplots(figsize=(4, 5))
ax.set_xlim(-3, 3)
ax.set_ylim(-4, 4)
ax.axis('off')

circle = patches.Circle((0, 0), 0.6, facecolor='whitesmoke', edgecolor='black', linewidth=1.5)
ax.add_patch(circle)

ax.annotate('', xy=(0, 2.5), xytext=(0, 0.6), arrowprops=dict(facecolor='black', width=2, headwidth=8))
ax.text(0.3, 2, 'R = bv', fontsize=14, fontweight='bold')

ax.annotate('', xy=(0, -2.5), xytext=(0, -0.6), arrowprops=dict(facecolor='black', width=2, headwidth=8))
ax.text(0.3, -2.2, 'P = mg', fontsize=14, fontweight='bold')

ax.annotate('', xy=(-2, 1), xytext=(-2, -1), arrowprops=dict(arrowstyle="->", color='gray'))
ax.text(-2.5, 0, 'y', color='gray', fontsize=12)

plt.title('DCL: Esfera en Rapidez Terminal', fontsize=14, fontweight='bold')
plt.tight_layout()
plt.show()
\end{lstlisting}

Aplicando la segunda ley de Newton en el eje vertical ($y$) durante la rapidez terminal, garantizamos que no hay aceleración:
$$\Sigma F_y = R - mg = 0$$
$$R = mg$$
Sustituyendo el modelo de la fuerza resistiva ($R = bv_T$) evaluado en la rapidez terminal:
$$bv_T = mg$$

\subsubsection*{Homogeneización del Sistema de Unidades}
Antes de manipular la ecuación algebraicamente, es obligatorio estandarizar todas las magnitudes al Sistema Internacional (SI) para evitar inconsistencias dimensionales en el cálculo de la constante.
Para la masa:
$$m = 2 \text{ g} = 0.002 \text{ kg}$$
Para la rapidez terminal:
$$v_T = 5 \text{ cm/s} = 0.05 \text{ m/s}$$

\subsubsection*{Cálculo de la Constante de Proporcionalidad}
Aislamos la constante resistiva $b$ de nuestra ecuación de equilibrio de fuerzas y sustituimos los valores estandarizados:
$$b = \frac{mg}{v_T}$$
$$b = \frac{(0.002)(10)}{0.05}$$
$$b = \frac{0.02}{0.05}$$
$$b = 0.4 \text{ kg/s}$$
Si se requiere expresar el resultado en las unidades originales del problema (gramos), multiplicamos por el factor de conversión ($1 \text{ kg} = 1000 \text{ g}$):
$$b = 0.4 \times 1000 = 400 \text{ g/s}$$

\subsubsection*{Conclusión}
La constante de resistividad $b$ del aceite es de $0.4 \text{ kg/s}$ (o equivalentemente $400 \text{ g/s}$).

\newpage

\subsection*{Enunciado}
Un paracaidista de 50 kg cae a través del aire. El coeficiente de arrastre para un cuerpo humano con extremidades extendidas es aproximadamente $D = 0.25 \text{ kg/m}$. Calcule su rapidez terminal utilizando el modelo $R = Dv^2$. Asuma la aceleración de la gravedad como $g = 10 \text{ m/s}^2$.

\subsection*{Solución}

\subsubsection*{Fundamento Físico de la Rapidez Terminal}
Siguiendo los principios de la física conceptual, cuando un paracaidista se lanza al vacío, la interacción gravitacional lo acelera hacia abajo. Sin embargo, a medida que su rapidez de caída aumenta, debe apartar las moléculas de aire de su camino con mayor intensidad, lo que genera una resistencia aerodinámica que lo empuja hacia arriba. Debido a las altas velocidades y a la turbulencia generada, esta fuerza de arrastre ($R$) es directamente proporcional al cuadrado de la rapidez ($v^2$). 
El concepto central de "rapidez terminal" se manifiesta cuando el paracaidista alcanza una velocidad exacta en la que la resistencia del aire crece lo suficiente como para igualar perfectamente la magnitud de su peso. En este estado de equilibrio dinámico, la fuerza neta se anula, la aceleración se vuelve cero y el paracaidista deja de ganar velocidad, cayendo con una rapidez máxima y constante.

\subsubsection*{Diagrama de Cuerpo Libre (Equilibrio Dinámico)}
Para representar gráficamente este estado de equilibrio, modelamos las fuerzas actuantes en el eje vertical con el siguiente script computacional.

\begin{lstlisting}[language=Python]
import matplotlib.pyplot as plt
import matplotlib.patches as patches

fig, ax = plt.subplots(figsize=(4, 5))
ax.set_xlim(-3, 3)
ax.set_ylim(-4, 4)
ax.axis('off')

rect = patches.Rectangle((-0.8, -0.6), 1.6, 1.2, facecolor='whitesmoke', edgecolor='black', linewidth=1.5)
ax.add_patch(rect)

ax.annotate('', xy=(0, 2.5), xytext=(0, 0.6), arrowprops=dict(facecolor='black', width=2, headwidth=8))
ax.text(0.3, 1.5, 'R = Dv^2', fontsize=14, fontweight='bold')

ax.annotate('', xy=(0, -2.5), xytext=(0, -0.6), arrowprops=dict(facecolor='black', width=2, headwidth=8))
ax.text(0.3, -1.8, 'P = mg', fontsize=14, fontweight='bold')

ax.annotate('', xy=(-2, 1), xytext=(-2, -1), arrowprops=dict(arrowstyle="->", color='gray'))
ax.text(-2.5, 0, 'y', color='gray', fontsize=12)

plt.title('DCL: Paracaidista en Rapidez Terminal', fontsize=12, fontweight='bold')
plt.tight_layout()
plt.show()
\end{lstlisting}

\subsubsection*{Aplicación de la Segunda Ley de Newton}
Dado que el paracaidista ya no acelera, aplicamos la condición de equilibrio traslacional en el eje vertical (eje $y$). La suma vectorial de las fuerzas debe ser estrictamente cero:
$$\Sigma F_y = R - mg = 0$$
$$R = mg$$
Sustituyendo la ecuación matemática del modelo de arrastre aerodinámico evaluado en el instante de la rapidez terminal ($v_T$):
$$Dv_T^2 = mg$$

\subsubsection*{Despeje y Cálculo Matemático}
Para aislar la variable cinemática de interés ($v_T$), dividimos ambos lados de la igualdad por el coeficiente de arrastre $D$ y aplicamos la raíz cuadrada para despejar el exponente cuadrático:
$$v_T^2 = \frac{mg}{D}$$
$$v_T = \sqrt{\frac{mg}{D}}$$
Procedemos a sustituir los valores numéricos proporcionados por el sistema ($m = 50 \text{ kg}$, $g = 10 \text{ m/s}^2$, $D = 0.25 \text{ kg/m}$):
$$v_T = \sqrt{\frac{(50)(10)}{0.25}}$$
$$v_T = \sqrt{\frac{500}{0.25}}$$
$$v_T = \sqrt{2000}$$
$$v_T \approx 44.72 \text{ m/s}$$
Para fines de comprensión intuitiva en la vida cotidiana, podemos transformar este resultado a kilómetros por hora aplicando el factor de conversión ($3.6$):
$$v_T \approx 161 \text{ km/h}$$

\subsubsection*{Conclusión}
La rapidez terminal que alcanza el paracaidista es de $44.72 \text{ m/s}$ (que equivale a aproximadamente $161 \text{ km/h}$).

\newpage

\subsection*{Enunciado}
Cuando empujas una pared, ¿la pared te empuja a ti? Si es así, ¿por qué no te mueves hacia atrás?

\subsection*{Solución}

\subsubsection*{El Principio de Interacción y la Tercera Ley de Newton}
De acuerdo con la física conceptual, las fuerzas nunca existen de forma aislada en el universo; siempre ocurren en pares que constituyen una interacción mutua entre dos objetos. La tercera ley de Newton establece que siempre que un objeto ejerce una fuerza sobre un segundo objeto, el segundo objeto ejerce una fuerza de igual magnitud y dirección opuesta sobre el primero. Por lo tanto, cuando tú empujas mecánicamente una pared con tus manos (fuerza de acción), es un hecho físico ineludible que la pared te empuja simultáneamente a ti con exactamente la misma cantidad de fuerza, pero en sentido contrario (fuerza de reacción). 

\subsubsection*{Análisis del Sistema y Fuerzas Externas}
Para comprender por qué no aceleras hacia atrás a pesar de recibir esta fuerza de la pared, debemos aislar tu cuerpo como el "sistema" de interés y analizar todas las fuerzas externas que actúan sobre ti. 
Mientras la pared te empuja hacia atrás, existe otro punto de contacto fundamental: tus pies sobre el suelo. Al intentar ser empujado hacia atrás, la suela de tus zapatos tiende a deslizarse sobre el piso. Esta tendencia al movimiento genera inmediatamente una fuerza de fricción estática ejercida por el suelo sobre tus zapatos, la cual apunta hacia adelante. 

\subsubsection*{Equilibrio Dinámico (Primera y Segunda Ley)}
Si no te mueves hacia atrás, significa que tu aceleración es cero. Según la segunda ley de Newton, esto requiere que la fuerza neta sobre tu cuerpo sea nula. Lo que ocurre físicamente es que la fuerza de fricción estática del suelo hacia adelante iguala y cancela exactamente la fuerza de empuje de la pared hacia atrás. 

\begin{lstlisting}[language=Python]
import matplotlib.pyplot as plt

fig, ax = plt.subplots(figsize=(6, 4))
ax.set_xlim(-4, 4)
ax.set_ylim(-1, 3)
ax.axis('off')

ax.plot([-4, 4], [0, 0], color='black', linewidth=2)
ax.plot([2, 2], [0, 3], color='gray', linewidth=4)

ax.plot([0, 0], [1, 2], color='blue', linewidth=3)
ax.plot([0, 1.8], [1.5, 1.5], color='blue', linewidth=2)
ax.plot([-0.5, 0.5], [0.1, 0.1], color='blue', linewidth=4)

ax.annotate('', xy=(-2, 1.5), xytext=(0, 1.5), arrowprops=dict(facecolor='red', width=3, headwidth=10))
ax.text(-1, 1.8, 'F_pared', color='red', fontsize=12, fontweight='bold', ha='center')

ax.annotate('', xy=(2, 0.3), xytext=(0, 0.3), arrowprops=dict(facecolor='green', width=3, headwidth=10))
ax.text(1, 0.6, 'F_friccion', color='green', fontsize=12, fontweight='bold', ha='center')

plt.title('Diagrama de Fuerzas sobre la Persona', fontsize=12, fontweight='bold')
plt.tight_layout()
plt.show()
\end{lstlisting}

\subsubsection*{Conclusión}
Sí, la pared te empuja hacia atrás con la misma fuerza con la que tú la empujas. Sin embargo, no te mueves hacia atrás porque la fuerza de fricción estática que el suelo ejerce sobre tus zapatos (hacia adelante) cancela el empuje de la pared, manteniendo tu cuerpo en un estado de equilibrio estático con una fuerza neta igual a cero.

\newpage

\subsection*{Enunciado}
Mencione al menos tres ejemplos de la vida cotidiana donde se presenta la tercera ley de Newton.

\subsection*{Solución}

\subsubsection*{Naturaleza de los Pares de Acción y Reacción}
La tercera ley de Newton requiere identificar claramente a los dos cuerpos que interactúan. La regla pedagógica fundamental de Hewitt para identificar estos pares es: si la acción es el objeto A ejerciendo fuerza sobre el objeto B, la reacción es obligatoriamente el objeto B ejerciendo fuerza sobre el objeto A. A continuación, se desglosan tres mecanismos cotidianos de locomoción que dependen enteramente de este principio.

\subsubsection*{Ejemplo 1: La Dinámica de Caminar}
Cuando una persona camina, no es la fuerza interna de sus músculos la que la desplaza directamente hacia adelante. La acción consiste en que el zapato empuja la superficie del suelo hacia atrás. Como reacción recíproca, el suelo empuja el zapato (y por ende a la persona) hacia adelante con una fuerza de fricción estática idéntica. Sin esta fuerza de reacción del suelo, caminar sería imposible.

\subsubsection*{Ejemplo 2: La Propulsión al Nadar}
Al nadar, los brazos y piernas del nadador desplazan físicamente el agua en dirección opuesta a la meta. La acción es el nadador empujando el volumen de agua hacia atrás. La reacción inmediata es que esa misma masa de agua empuja al nadador hacia adelante. El movimiento en el fluido se logra exclusivamente aprovechando esta fuerza de reacción.

\subsubsection*{Ejemplo 3: El Vuelo de las Aves y Aeronaves}
El vuelo sostenido ilustra el mismo principio tridimensional. La acción consiste en que las alas del pájaro (o las aspas de un helicóptero) empujan las moléculas de aire hacia abajo con gran fuerza. En respuesta, la tercera ley dicta que el aire reaccione empujando las alas hacia arriba. Esta fuerza de reacción ascendente es lo que conocemos como sustentación, la cual compensa el peso de la criatura.

\subsubsection*{Conclusión}
Tres ejemplos fundamentales son: caminar (el pie empuja el suelo hacia atrás y el suelo empuja a la persona hacia adelante), nadar (la mano empuja el agua hacia atrás y el agua empuja al nadador hacia adelante) y el vuelo de un ave (las alas empujan el aire hacia abajo y el aire empuja al ave hacia arriba).

\newpage

\subsection*{Enunciado}
Durante el choque entre un insecto y un autobús, ¿quién experimenta una fuerza mayor? ¿y aceleración mayor?

\subsection*{Solución}

\subsubsection*{Análisis de la Magnitud de la Fuerza de Impacto}
Es un instinto común suponer que el objeto más grande y pesado ejerce una fuerza mayor. Sin embargo, la física conceptual requiere que evaluemos el evento como una sola interacción mutua. El impacto no es una fuerza que el autobús "posee" y le "entrega" al insecto; es un evento compartido entre las dos masas. De acuerdo estricto con la tercera ley de Newton, la fuerza que el autobús ejerce sobre el insecto (acción) es exactamente igual en magnitud y opuesta en dirección a la fuerza que el insecto ejerce sobre el autobús (reacción). Ambos experimentan numéricamente la misma cantidad de fuerza durante el milisegundo que dura la colisión.

\subsubsection*{Análisis de la Respuesta Cinemática (Aceleración)}
La diferencia dramática en los resultados del choque no proviene de una diferencia en las fuerzas, sino de una diferencia extrema en las inercias (masas) de los objetos, analizada mediante la segunda ley de Newton. Esta ley dicta que la aceleración es directamente proporcional a la fuerza neta e inversamente proporcional a la masa ($a = F / m$). 

Para el insecto, que posee una masa minúscula ($m$), dividir la fuerza del impacto ($F$) entre un número tan cercano a cero produce una aceleración de desaceleración gigantesca. Es este cambio brusco de velocidad lo que resulta letal para su estructura.
Para el autobús, que posee una masa colosal ($M$), dividir la misma fuerza ($F$) entre un número inmenso resulta en una aceleración casi indetectable. Su estado de movimiento apenas es perturbado.

\begin{lstlisting}[language=Python]
import matplotlib.pyplot as plt

fig, ax = plt.subplots(figsize=(7, 4))
ax.axis('off')

ax.text(0.2, 0.8, 'Insecto (m pequeña)', fontsize=12, fontweight='bold')
ax.annotate('', xy=(0.8, 0.7), xytext=(0.1, 0.7), arrowprops=dict(facecolor='red', width=3, headwidth=10))
ax.text(0.4, 0.55, 'a = F / m (Aceleracion Enorme)', fontsize=11)

ax.text(0.2, 0.3, 'Autobus (M enorme)', fontsize=12, fontweight='bold')
ax.annotate('', xy=(0.4, 0.2), xytext=(0.1, 0.2), arrowprops=dict(facecolor='red', width=3, headwidth=10))
ax.text(0.4, 0.05, 'A = F / M (Aceleracion Minima)', fontsize=11)

plt.title('Comparacion de Aceleraciones bajo Fuerzas Iguales', fontsize=14, fontweight='bold')
plt.tight_layout()
plt.show()
\end{lstlisting}

\subsubsection*{Conclusión}
Ambos cuerpos, el insecto y el autobús, experimentan exactamente la misma magnitud de fuerza durante el choque, cumpliendo la tercera ley de Newton. Sin embargo, debido a que el insecto tiene una masa infinitesimalmente menor en comparación con el autobús, experimenta una aceleración infinitamente mayor de acuerdo con la segunda ley de Newton.

\newpage

\subsection*{Enunciado}
Dos bloques de masas $m_1$ y $m_2$, con $m_1 > m_2$, se colocan en contacto mutuo sobre una superficie horizontal sin fricción. Una fuerza horizontal constante $\vec{F}$ se aplica a $m_1$ como se muestra en la figura. Determine la magnitud de la aceleración del sistema.

\begin{center}
\begin{tikzpicture}
\fill[gray!30] (-2,-0.2) rectangle (4,0);
\draw[thick] (-2,0) -- (4,0);
\draw[thick, fill=orange!40] (0,0) rectangle (1.2,1);
\node at (0.6,0.5) {$m_1$};
\draw[thick, fill=cyan!40] (1.2,0) rectangle (2,0.8);
\node at (1.6,0.4) {$m_2$};
\draw[->, very thick, cyan] (-1.5,0.5) -- (-0.1,0.5) node[midway, above, text=black] {$\vec{F}$};
\end{tikzpicture}
\end{center}

\subsection*{Solución}

\subsubsection*{Conceptualización del Sistema Macroscópico}
De acuerdo con los principios de la física conceptual de Hewitt, cuando múltiples objetos se mueven juntos en la misma dirección y al mismo ritmo debido a que están en contacto directo, es estratégicamente útil analizarlos inicialmente como si fueran un solo objeto sólido. Dado que los bloques $m_1$ y $m_2$ son empujados juntos y la superficie no presenta fricción, ambos compartirán idéntica aceleración. Por lo tanto, podemos definir nuestro "sistema" como la combinación de ambas masas, teniendo una masa total equivalente a $M = m_1 + m_2$. Al definir el sistema de esta manera abarcadora, las fuerzas internas (como el empuje entre los bloques) se cancelan y no afectan el movimiento global.

\subsubsection*{Aplicación de la Segunda Ley de Newton}
La única fuerza externa neta que actúa horizontalmente sobre nuestro sistema combinado es la fuerza aplicada $\vec{F}$. Aplicando la segunda ley de Newton en el eje horizontal ($\Sigma F_x = M \cdot a_x$), igualamos esta fuerza a la masa total del sistema multiplicada por su aceleración conjunta.
$$F = (m_1 + m_2)a$$
Para encontrar la aceleración de los bloques, simplemente realizamos un despeje algebraico dividiendo la fuerza aplicada por la masa total.
$$a = \frac{F}{m_1 + m_2}$$

\subsubsection*{Conclusión}
La magnitud de la aceleración del sistema combinado es $a = \frac{F}{m_1 + m_2}$.

\newpage

\subsection*{Enunciado}
Dos bloques de masas $m_1$ y $m_2$, con $m_1 > m_2$, se colocan en contacto mutuo sobre una superficie horizontal sin fricción. Una fuerza horizontal constante $\vec{F}$ se aplica a $m_1$ como se muestra en la figura. Determine la magnitud de la fuerza de contacto entre los dos bloques.

\begin{center}
\begin{tikzpicture}
\fill[gray!30] (-2,-0.2) rectangle (4,0);
\draw[thick] (-2,0) -- (4,0);
\draw[thick, fill=orange!40] (0,0) rectangle (1.2,1);
\node at (0.6,0.5) {$m_1$};
\draw[thick, fill=cyan!40] (1.2,0) rectangle (2,0.8);
\node at (1.6,0.4) {$m_2$};
\draw[->, very thick, cyan] (-1.5,0.5) -- (-0.1,0.5) node[midway, above, text=black] {$\vec{F}$};
\end{tikzpicture}
\end{center}

\subsection*{Solución}

\subsubsection*{Tercera Ley de Newton y Aislamiento de Subsistemas}
Para determinar la fuerza de contacto entre las superficies de los bloques, debemos abandonar la visión del "sistema combinado" y enfocarnos en las interacciones individuales. Por la tercera ley de Newton (acción y reacción), sabemos que el bloque $m_1$ empuja al bloque $m_2$ hacia la derecha con una fuerza de contacto $F_c$. Inmediatamente, el bloque $m_2$ reacciona empujando al bloque $m_1$ hacia la izquierda con una fuerza de idéntica magnitud $F_c$. Para hallar matemáticamente esta magnitud, es pedagógicamente más sencillo aislar el bloque $m_2$ y establecerlo como nuestro nuevo sistema de interés.

\subsubsection*{Representación Gráfica de Diagramas de Cuerpo Libre}
El siguiente modelado ilustra las diferencias en las fuerzas horizontales que gobiernan a cada bloque por separado. Note que sobre $m_2$, la única fuerza horizontal que produce aceleración es la fuerza de contacto $F_c$.

\begin{lstlisting}[language=Python]
import matplotlib.pyplot as plt

fig, (ax1, ax2) = plt.subplots(1, 2, figsize=(8, 4))

ax1.set_xlim(-3, 3)
ax1.set_ylim(-3, 3)
ax1.axis('off')
ax1.add_patch(plt.Rectangle((-0.8, -0.6), 1.6, 1.2, fill=True, color='bisque', ec='black'))
ax1.text(0, 0, 'm1', ha='center', va='center', fontsize=14)
ax1.annotate('', xy=(0, 2), xytext=(0, 0.6), arrowprops=dict(facecolor='black', width=2, headwidth=8))
ax1.text(0.2, 1.5, 'N1', fontsize=12, fontweight='bold')
ax1.annotate('', xy=(0, -2), xytext=(0, -0.6), arrowprops=dict(facecolor='black', width=2, headwidth=8))
ax1.text(0.2, -1.8, 'P1', fontsize=12, fontweight='bold')
ax1.annotate('', xy=(2, 0), xytext=(0.8, 0), arrowprops=dict(facecolor='blue', width=2, headwidth=8))
ax1.text(1.2, 0.3, 'F', color='blue', fontsize=12, fontweight='bold')
ax1.annotate('', xy=(-2, 0), xytext=(-0.8, 0), arrowprops=dict(facecolor='red', width=2, headwidth=8))
ax1.text(-1.8, 0.3, 'Fc', color='red', fontsize=12, fontweight='bold')
ax1.set_title('DCL: Bloque m1', fontweight='bold')

ax2.set_xlim(-3, 3)
ax2.set_ylim(-3, 3)
ax2.axis('off')
ax2.add_patch(plt.Rectangle((-0.6, -0.5), 1.2, 1.0, fill=True, color='lightcyan', ec='black'))
ax2.text(0, 0, 'm2', ha='center', va='center', fontsize=14)
ax2.annotate('', xy=(0, 2), xytext=(0, 0.5), arrowprops=dict(facecolor='black', width=2, headwidth=8))
ax2.text(0.2, 1.5, 'N2', fontsize=12, fontweight='bold')
ax2.annotate('', xy=(0, -2), xytext=(0, -0.5), arrowprops=dict(facecolor='black', width=2, headwidth=8))
ax2.text(0.2, -1.8, 'P2', fontsize=12, fontweight='bold')
ax2.annotate('', xy=(2, 0), xytext=(0.6, 0), arrowprops=dict(facecolor='red', width=2, headwidth=8))
ax2.text(1.2, 0.3, 'Fc', color='red', fontsize=12, fontweight='bold')
ax2.set_title('DCL: Bloque m2', fontweight='bold')

plt.tight_layout()
plt.show()
\end{lstlisting}

\subsubsection*{Cálculo de la Fuerza de Contacto}
Concéntrate exclusivamente en el diagrama del bloque $m_2$. La sumatoria de fuerzas en su eje horizontal es extremadamente simple, pues carece de fricción y no recibe el empuje externo directo $\vec{F}$. La segunda ley de Newton establece que la fuerza neta es el producto de su propia masa y su aceleración ($\Sigma F_{x,2} = m_2 \cdot a$).
$$F_c = m_2 \cdot a$$
Dado que la aceleración ($a$) es una constante compartida en el sistema calculada en la primera parte del problema, la sustituimos directamente en la ecuación.
$$F_c = m_2 \left( \frac{F}{m_1 + m_2} \right)$$
Reordenando los términos para una representación matemática más elegante, observamos que la fuerza de contacto es una fracción pura de la fuerza total aplicada.
$$F_c = \left( \frac{m_2}{m_1 + m_2} \right) F$$

\subsubsection*{Conclusión}
La magnitud de la fuerza de contacto que empuja a los bloques entre sí es $F_c = \frac{m_2}{m_1 + m_2}F$.

\newpage

\subsection*{Enunciado}
Si la Tierra ejerce una fuerza gravitacional de magnitud $mg$ sobre un proyectil en movimiento vertical de caída libre, hacia abajo, ¿cuál es la reacción a esta fuerza?

a) Una fuerza de magnitud cero, dado que la masa de la Tierra es inmensa.

b) La fuerza normal que el aire ejerce hacia arriba sobre el proyectil.

c) La fuerza gravitacional que ejerce el proyectil sobre la Tierra.

d) La fuerza de inercia que mantiene al proyectil en movimiento.

e) La aceleración del proyectil hacia el centro de la Tierra.

\subsection*{Solución}

\subsubsection*{Definición de un Par de Acción y Reacción}
En la física conceptual, la tercera ley de Newton establece una regla rigurosa y simétrica para identificar las interacciones mutuas: si un objeto A ejerce una fuerza sobre un objeto B (acción), entonces el objeto B ejerce simultáneamente una fuerza de igual magnitud y dirección opuesta sobre el objeto A (reacción). Es fundamental entender que el par de fuerzas actúa exclusivamente sobre estos dos cuerpos que interactúan; no involucra a un tercer cuerpo ni a otros fenómenos físicos.

\subsubsection*{Análisis de la Interacción Gravitacional}
Aplicando esta regla al problema, definimos a los dos actores de la interacción. El objeto A es la Tierra y el objeto B es el proyectil. 
La fuerza de "acción" se describe como: La Tierra (A) tira del proyectil (B) hacia abajo con una fuerza gravitacional de magnitud $mg$.
Por estricta simetría matemática y conceptual, la fuerza de "reacción" debe ser: El proyectil (B) tira de la Tierra (A) hacia arriba con una fuerza gravitacional exactamente de la misma magnitud $mg$.
Es un error común creer que, por ser la Tierra inmensamente más masiva, no sufre una fuerza de reacción (opción a). La fuerza existe y es de la misma magnitud; sin embargo, debido a la segunda ley de Newton ($a = \frac{F}{m}$), aplicar una fuerza tan pequeña ($mg$) a la inmensa masa planetaria produce una aceleración infinitesimalmente imperceptible, por lo que no observamos a la Tierra moverse hacia el proyectil.

\subsubsection*{Conclusión}
La opción correcta es la c. La reacción a la fuerza gravitacional de la Tierra sobre el proyectil es la fuerza gravitacional que ejerce el proyectil sobre la Tierra.

\newpage
\subsection*{Enunciado}
Un camión de gran tonelaje choca frontalmente con un automóvil pequeño. Durante el impacto, ¿cuál de las siguientes opciones describe correctamente las fuerzas según la tercera ley de Newton?

a) Las fuerzas son distintas porque el automóvil experimenta una aceleración mucho mayor.

b) El automóvil ejerce una fuerza mayor para intentar detener la inercia del camión.

c) El camión ejerce una fuerza mayor debido a que posee una mayor masa.

d) La fuerza ejercida por el camión es igual en magnitud a la fuerza ejercida por el automóvil.

e) No existen fuerzas de reacción si los objetos quedan unidos después del choque.

\subsection*{Solución}

\subsubsection*{El Principio de Interacción Mutua en Colisiones}
Un choque físico no es una transferencia unilateral de fuerza, sino una interacción mutua y compartida entre dos cuerpos. De acuerdo con la tercera ley de Newton, toda fuerza es parte de una interacción simétrica. En el instante exacto en que los parachoques entran en contacto, el camión empuja al automóvil y, simultáneamente, el automóvil empuja al camión. Estas dos fuerzas constituyen un par de acción y reacción inquebrantable.

\subsubsection*{Independencia entre Fuerza y Aceleración}
Las opciones a, b y c representan conceptos erróneos muy arraigados. Se suele confundir la "fuerza del impacto" con el "daño" o la "aceleración" resultante. La tercera ley garantiza que la fuerza de impacto es idéntica para ambos vehículos ($F_{\text{camión sobre auto}} = -F_{\text{auto sobre camión}}$). La razón por la cual el automóvil pequeño queda destrozado y sale despedido hacia atrás no es porque reciba más fuerza, sino porque posee mucha menos masa. Según la segunda ley de Newton ($a = \frac{F}{m}$), esa misma magnitud de fuerza compartida produce una aceleración (y deformación estructural) catastrófica en la masa pequeña del auto, pero apenas perturba el estado de movimiento de la enorme masa del camión.

\subsubsection*{Representación Gráfica del Choque}
Para visualizar pedagógicamente la diferencia entre las inercias (masas) y la igualdad innegable de las fuerzas de interacción, se presenta la siguiente simulación del instante del impacto.

\begin{lstlisting}[language=Python]
import matplotlib.pyplot as plt
import matplotlib.patches as patches

fig, ax = plt.subplots(figsize=(8, 4))
ax.set_xlim(-5, 5)
ax.set_ylim(-2, 3)
ax.axis('off')

camion = patches.Rectangle((-4, -1), 3.5, 2.5, facecolor='lightgray', edgecolor='black', linewidth=2)
ax.add_patch(camion)
ax.text(-2.25, 0.25, 'Camion\n(Gran Masa)', ha='center', va='center', fontsize=12, fontweight='bold')

auto = patches.Rectangle((0.5, -1), 2, 1.2, facecolor='lightblue', edgecolor='black', linewidth=2)
ax.add_patch(auto)
ax.text(1.5, -0.4, 'Auto\n(Poca Masa)', ha='center', va='center', fontsize=11, fontweight='bold')

ax.annotate('', xy=(0.5, 0), xytext=(-0.5, 0), arrowprops=dict(facecolor='red', width=4, headwidth=12))
ax.text(0, 0.5, 'F (Auto sobre Camion)', color='red', fontsize=10, fontweight='bold', ha='center')

ax.annotate('', xy=(-0.5, -0.5), xytext=(0.5, -0.5), arrowprops=dict(facecolor='blue', width=4, headwidth=12))
ax.text(0, -1, 'F (Camion sobre Auto)', color='blue', fontsize=10, fontweight='bold', ha='center')

plt.title('Interaccion Mutua en Choque: Fuerzas Iguales y Opuestas', fontsize=14, fontweight='bold')
plt.tight_layout()
plt.show()
\end{lstlisting}

\subsubsection*{Conclusión}
La opción correcta es la d. Independientemente de la masa o de los daños sufridos, la fuerza ejercida por el camión es estrictamente igual en magnitud a la fuerza ejercida por el automóvil.

\newpage
\subsection*{Enunciado}
Se afirma que un balón de fútbol no acelera cuando dos personas lo patean simultáneamente con fuerzas de igual magnitud y opuestas. ¿Por qué estas dos fuerzas (de las patadas) no forman un par acción-reacción?

a) Porque ambas fuerzas actúan sobre el mismo cuerpo.

b) Porque las patadas no son del mismo tipo de interacción física.

c) Porque se requieren al menos tres objetos para formar un par de acción-reacción.

d) Porque las fuerzas de reacción solo existen en objetos inertes.

e) Porque la fuerza neta sobre el balón es cero.

\subsection*{Solución}

\subsubsection*{Naturaleza de los Pares de la Tercera Ley de Newton}
El pilar fundamental para identificar un par de fuerzas de acción y reacción es reconocer que jamás pueden sumarse o anularse entre sí matemáticamente para producir equilibrio. Esto se debe a una regla inquebrantable de la tercera ley de Newton: las fuerzas de acción y reacción actúan siempre sobre cuerpos distintos (A sobre B, y B sobre A).

\subsubsection*{Equilibrio vs. Acción-Reacción}
En el escenario descrito, la Persona 1 aplica una fuerza sobre el balón, y la Persona 2 aplica una fuerza opuesta también sobre el balón. Como ambas fuerzas actúan sobre el idéntico objeto físico (el balón), estas pueden sumarse vectorialmente. Al ser iguales y opuestas, la suma da como resultado una fuerza neta de cero ($\Sigma F = 0$). Este es un caso clásico de la primera y segunda ley de Newton, que genera un estado de equilibrio traslacional donde el balón no acelera.
Sin embargo, estas dos fuerzas no constituyen un par de acción-reacción. El verdadero par para la patada de la Persona 1 es la fuerza de reacción del balón aplastando o empujando el pie de esa misma persona hacia atrás. 

\subsubsection*{Conclusión}
La opción correcta es la a. Dos fuerzas que se cancelan en un sistema en equilibrio no forman un par acción-reacción porque actúan simultáneamente sobre el mismo cuerpo (el balón).

\newpage
\subsection*{Enunciado}
De acuerdo con el ejemplo del helicóptero en los textos proporcionados, ¿cómo se genera la fuerza de sustentación?

a) Las hélices empujan el aire hacia abajo y, en reacción, el aire empuja el helicóptero hacia arriba.

b) El motor del helicóptero anula la fuerza de gravedad internamente.

c) La fuerza de sustentación es una propiedad intrínseca del giro de las hélices, sin necesidad de interactuar con el aire.

d) La fricción del aire contra el fuselaje empuja al helicóptero hacia arriba.

e) Las hélices crean un vacío sobre el helicóptero que lo succiona hacia arriba.

\subsection*{Solución}

\subsubsection*{Interacción con Fluidos y Mecánica de Vuelo}
Para comprender cómo un objeto pesado logra vencer la gravedad y elevarse en el aire, es necesario recurrir a la tercera ley de Newton aplicada a la mecánica de fluidos. Ningún objeto posee la capacidad aislada de "elevarse a sí mismo" sin interactuar con su entorno material. El vuelo se logra estrictamente mediante la manipulación de una gran masa del fluido circundante (en este caso, el aire atmosférico).

\subsubsection*{Mecánica de Sustentación en Helicópteros}
Las aspas del rotor de un helicóptero no son planas; están diseñadas con un perfil aerodinámico y un ángulo de ataque específico. Cuando el motor las hace girar a altas velocidades, las aspas golpean físicamente las moléculas de aire y las aceleran con gran fuerza hacia el suelo. Esta interacción constituye la fuerza de "acción" (hélices empujando el aire hacia abajo).
En concordancia inmediata con la tercera ley de Newton, la respuesta ineludible de la naturaleza es que esa misma masa de aire reaccione empujando las aspas hacia arriba con una fuerza de idéntica magnitud. Esta fuerza reactiva ascendente, transmitida de las aspas al resto de la aeronave, es lo que conocemos técnicamente como fuerza de sustentación. Cuando esta reacción supera el peso total del aparato, el helicóptero asciende.

\subsubsection*{Conclusión}
La opción correcta es la a. La sustentación es la consecuencia directa de la tercera ley de Newton: las hélices realizan la acción de empujar el aire hacia abajo, y como reacción, el fluido (aire) empuja al helicóptero hacia arriba.

\newpage

\subsection*{Enunciado}
De acuerdo con la tercera ley de Newton, una fuerza individual aislada:

a) puede existir si el objeto tiene mucha masa.

b) solo existe en marcos de referencia no inerciales.

c) no puede existir, ya que las fuerzas son siempre parte de una interacción mutua.

d) solo existe cuando un objeto está en equilibrio estático.

e) es posible siempre que no haya fricción en el medio.

\subsection*{Solución}

\subsubsection*{La Naturaleza de las Fuerzas como Interacciones}
En el marco de la física conceptual propuesto por Hewitt, es fundamental cambiar la perspectiva coloquial de lo que es una fuerza. Una fuerza no es una "cosa" que un objeto posea en su interior, de la misma manera que posee masa o volumen. Por el contrario, una fuerza es estrictamente un empuje o un tirón que ocurre como resultado de una interacción entre dos entidades físicas distintas. 

\subsubsection*{La Imposibilidad de una Fuerza Aislada}
Puesto que una fuerza es el producto de una interacción, lógicamente requiere la participación de al menos dos cuerpos: un cuerpo que ejerce la acción y un cuerpo que recibe dicha acción y reacciona. Hewitt utiliza una analogía insuperable para ilustrar este concepto: es físicamente imposible que toques algo sin ser tocado al mismo tiempo y con la misma intensidad. Si golpeas una pared, la pared golpea tu mano. Por lo tanto, el concepto de una "fuerza individual aislada" actuando sola en el universo viola el principio fundamental de la tercera ley de Newton. Las fuerzas siempre ocurren en pares inseparables (acción y reacción).

\subsubsection*{Representación Gráfica de la Interacción Mutua}
Para visualizar esta dependencia estricta, el siguiente modelado computacional muestra cómo cualquier contacto o interacción a distancia genera instantáneamente dos vectores de fuerza de igual magnitud sobre cuerpos distintos.

\begin{lstlisting}[language=Python]
import matplotlib.pyplot as plt

fig, ax = plt.subplots(figsize=(6, 3))
ax.set_xlim(-3, 3)
ax.set_ylim(-2, 2)
ax.axis('off')

circleA = plt.Circle((-1, 0), 0.8, facecolor='lightblue', edgecolor='black', linewidth=2)
circleB = plt.Circle((1, 0), 0.8, facecolor='lightgreen', edgecolor='black', linewidth=2)
ax.add_patch(circleA)
ax.add_patch(circleB)

ax.text(-1, 0, 'Objeto A', ha='center', va='center', fontsize=12, fontweight='bold')
ax.text(1, 0, 'Objeto B', ha='center', va='center', fontsize=12, fontweight='bold')

ax.annotate('', xy=(0.2, 1.2), xytext=(-1, 1.2), arrowprops=dict(facecolor='red', width=3, headwidth=10))
ax.text(-0.4, 1.5, 'Fuerza de A sobre B', color='red', fontsize=10, fontweight='bold', ha='center')

ax.annotate('', xy=(-0.2, -1.2), xytext=(1, -1.2), arrowprops=dict(facecolor='blue', width=3, headwidth=10))
ax.text(0.4, -1.6, 'Fuerza de B sobre A', color='blue', fontsize=10, fontweight='bold', ha='center')

plt.title('Interaccion Mutua: Las Fuerzas Siempre Vienen en Pares', fontsize=12, fontweight='bold')
plt.tight_layout()
plt.show()
\end{lstlisting}

\subsubsection*{Conclusión}
La opción correcta es la c. Una fuerza individual aislada no puede existir en la naturaleza, ya que, por definición, todas las fuerzas son el resultado de una interacción mutua entre dos objetos.

\newpage

\subsection*{Enunciado}
¿Por qué las fuerzas de acción y reacción nunca se cancelan entre sí?

a) Porque son fuerzas de diferentes tipos, como gravitacional y normal.

b) Porque actúan en tiempos diferentes durante la interacción.

c) Porque la magnitud de una es siempre ligeramente superior a la otra.

d) Porque actúan sobre objetos diferentes, por lo que no pueden sumarse para dar una fuerza neta sobre un solo cuerpo.

e) Porque solo se cancelan cuando el sistema está en el vacío absoluto.

\subsection*{Solución}

\subsubsection*{La Aparente Paradoja del Movimiento}
Desde la perspectiva de la física conceptual, la tercera ley de Newton a menudo genera una confusión clásica conocida como la "paradoja del movimiento". El razonamiento erróneo postula que, si toda fuerza de acción genera una fuerza de reacción exactamente de la misma magnitud pero en dirección opuesta, ambas fuerzas deberían sumarse algebraicamente y dar como resultado cero. Si la fuerza neta siempre fuera cero, ningún objeto en el universo podría acelerar o moverse. Para resolver esta paradoja, es fundamental comprender cómo se aplican las leyes de Newton a sistemas definidos.

\subsubsection*{El Principio de Aislamiento del Sistema}
La segunda ley de Newton ($\Sigma F = ma$) es la herramienta que utilizamos para predecir si un objeto va a acelerar. La regla de oro para aplicar esta ley es que, para calcular la aceleración de un cuerpo específico, única y exclusivamente debemos sumar las fuerzas externas que actúan sobre ese cuerpo en particular. Las fuerzas que ese cuerpo ejerce sobre otros elementos del entorno son irrelevantes para su propio estado de movimiento.

\subsubsection*{Separación de los Pares de Interacción}
En cualquier interacción mutua, existen dos cuerpos participantes. Si el Cuerpo A empuja al Cuerpo B (acción), esa fuerza actúa físicamente sobre la estructura del Cuerpo B. Como respuesta, el Cuerpo B empuja al Cuerpo A (reacción), y esta fuerza actúa físicamente sobre la estructura del Cuerpo A. Dado que cada vector de fuerza se dibuja en un Diagrama de Cuerpo Libre completamente distinto, es matemáticamente incorrecto e imposible sumarlos entre sí para buscar una cancelación. Las fuerzas solo se cancelan (produciendo equilibrio) cuando actúan simultáneamente en direcciones opuestas sobre el mismo objeto.

\subsubsection*{Representación Gráfica de la Separación Vectorial}
Para visualizar pedagógicamente por qué estas fuerzas no pueden anularse mutuamente, el siguiente modelado computacional grafica la interacción, demostrando que cada vector fuerza le pertenece a un objeto diferente.

\begin{lstlisting}[language=Python]
import matplotlib.pyplot as plt

fig, ax = plt.subplots(figsize=(7, 3.5))
ax.set_xlim(-4, 4)
ax.set_ylim(-2, 2)
ax.axis('off')

ax.add_patch(plt.Rectangle((-2.5, -0.6), 1.2, 1.2, facecolor='lightblue', edgecolor='black', linewidth=1.5))
ax.text(-1.9, 0, 'Cuerpo\nA', ha='center', va='center', fontsize=12, fontweight='bold')

ax.add_patch(plt.Rectangle((1.3, -0.6), 1.2, 1.2, facecolor='lightcoral', edgecolor='black', linewidth=1.5))
ax.text(1.9, 0, 'Cuerpo\nB', ha='center', va='center', fontsize=12, fontweight='bold')

ax.annotate('', xy=(1.3, 0.3), xytext=(-1.3, 0.3), arrowprops=dict(facecolor='red', width=2, headwidth=8))
ax.text(0, 0.6, 'Accion (Fuerza sobre B)', color='red', fontsize=11, fontweight='bold', ha='center')

ax.annotate('', xy=(-1.3, -0.3), xytext=(1.3, -0.3), arrowprops=dict(facecolor='blue', width=2, headwidth=8))
ax.text(0, -0.7, 'Reaccion (Fuerza sobre A)', color='blue', fontsize=11, fontweight='bold', ha='center')

plt.title('Separacion del Par de Fuerzas de la Tercera Ley', fontsize=12, fontweight='bold')
plt.tight_layout()
plt.show()
\end{lstlisting}

\subsubsection*{Conclusión}
La opción correcta es la d. Aunque las fuerzas de acción y reacción son perfectamente iguales y opuestas en cualquier interacción, nunca se anulan ni se cancelan entre sí porque actúan sobre objetos diferentes. Para obtener la fuerza neta sobre un solo cuerpo, no se pueden sumar fuerzas que se están aplicando a cuerpos distintos.

\newpage

\subsection*{Enunciado}
Una mosca choca con el parabrisas de un autobús que se mueve rápidamente. Durante el impacto:

a) la mosca experimenta una fuerza de impacto mucho mayor que el autobús.

b) el autobús experimenta una fuerza mayor debido a su gran energía cinética.

c) ambos experimentan la misma magnitud de fuerza de impacto.

d) la fuerza sobre la mosca depende solo de su propia masa.

e) no hay interacción de fuerzas si el choque es perfectamente inelástico.

\subsection*{Solución}

\subsubsection*{El Principio de Interacción Mutua}
Para analizar este evento desde la perspectiva de la física conceptual, debemos recordar que una fuerza no es una propiedad aislada que un objeto "posee" y "arroja" contra otro. Una fuerza es, por definición, una interacción mutua entre dos cuerpos. La tercera ley de Newton establece inequívocamente que toda fuerza de acción está acompañada por una fuerza de reacción de idéntica magnitud y dirección opuesta. En el instante del choque, el autobús empuja a la mosca (acción) y, de forma simultánea, la mosca empuja al autobús (reacción). 

\subsubsection*{Desmitificación del Tamaño y la Fuerza}
El sentido común a menudo nos engaña haciéndonos creer que el objeto más masivo, más grande o más rápido debe ejercer una fuerza "mayor". Sin embargo, esto es físicamente imposible. La fuerza de impacto es exactamente la misma para ambos. La razón por la que la mosca queda aplastada mientras que el autobús no sufre daños perceptibles no se debe a una diferencia en la magnitud de la fuerza, sino a una diferencia extrema en sus masas. 
Aplicando la segunda ley de Newton ($a = \frac{F}{m}$), la misma fuerza $F$ aplicada a la diminuta masa de la mosca produce una aceleración (desaceleración) letal e inmensa. Esa misma fuerza $F$ aplicada a la masa colosal del autobús produce una aceleración tan microscópica que no altera su estado de movimiento.

\subsubsection*{Representación Gráfica de la Interacción Simétrica}
Para ilustrar pedagógicamente que la asimetría de los tamaños no afecta la simetría de las fuerzas, se presenta el siguiente modelado computacional. Obsérvese que los vectores de fuerza son idénticos en longitud, independientemente de la masa del objeto que los recibe.

\begin{lstlisting}[language=Python]
import matplotlib.pyplot as plt
import matplotlib.patches as patches

fig, ax = plt.subplots(figsize=(8, 3))
ax.set_xlim(-5, 5)
ax.set_ylim(-2, 2)
ax.axis('off')

autobus = patches.Rectangle((-4, -1), 3.5, 2, facecolor='lightgray', edgecolor='black', linewidth=2)
ax.add_patch(autobus)
ax.text(-2.25, 0, 'Autobus\n(Gran Masa)', ha='center', va='center', fontsize=12, fontweight='bold')

mosca = patches.Circle((2, 0), 0.2, facecolor='black')
ax.add_patch(mosca)
ax.text(2, -0.6, 'Mosca\n(Poca Masa)', ha='center', va='center', fontsize=10, fontweight='bold')

ax.annotate('', xy=(1.5, 0.2), xytext=(-0.5, 0.2), arrowprops=dict(facecolor='red', width=3, headwidth=10))
ax.text(0.5, 0.5, 'Fuerza sobre la Mosca', color='red', fontsize=10, fontweight='bold', ha='center')

ax.annotate('', xy=(-0.5, -0.2), xytext=(1.5, -0.2), arrowprops=dict(facecolor='blue', width=3, headwidth=10))
ax.text(0.5, -0.6, 'Fuerza sobre el Autobus', color='blue', fontsize=10, fontweight='bold', ha='center')

plt.title('Tercera Ley de Newton: Interaccion Estrictamente Simetrica', fontsize=12, fontweight='bold')
plt.tight_layout()
plt.show()
\end{lstlisting}

\subsubsection*{Conclusión}
La opción correcta es la c. Independientemente de la enorme diferencia en masa y velocidad, la interacción es recíproca. Ambos cuerpos, el autobús y la mosca, experimentan exactamente la misma magnitud de fuerza de impacto durante la colisión.

\newpage

\subsection*{Enunciado}
En una "guerra de jalar la soga", el equipo ganador es aquel que:

a) tira de la soga con una fuerza mayor a la que el otro equipo tira de ella.

b) empuja con más fuerza contra el suelo, provocando una mayor fuerza de reacción externa sobre ellos.

c) tiene más masa, eliminando la necesidad de aplicar fuerzas externas.

d) logra que la tensión en su extremo de la soga sea mayor que en el otro.

e) reduce su propia inercia para acelerar más rápido hacia atrás.

\subsection*{Solución}

\subsubsection*{Análisis de las Fuerzas Internas y la Tensión}
Para resolver este problema, debemos basarnos en las leyes del movimiento de Newton, específicamente en la Tercera Ley (acción y reacción), tal como se aborda en la Física Conceptual. Un error común es pensar que el equipo ganador "tira con más fuerza de la soga". Sin embargo, la soga actúa únicamente como un transmisor de fuerza. Si asumimos que la masa de la soga es despreciable, la tensión a lo largo de toda la cuerda es exactamente la misma en cualquier punto. Esto significa que la fuerza con la que el equipo A tira del equipo B es idéntica en magnitud, pero opuesta en dirección, a la fuerza con la que el equipo B tira del equipo A. Por lo tanto, las opciones que sugieren diferencias en la tensión o en la fuerza de tracción directa sobre la cuerda son incorrectas, ya que estas representan fuerzas internas del sistema compuesto por ambos equipos y la soga.

\subsubsection*{El Rol de las Fuerzas Externas}
Según la Segunda Ley de Newton, para que un cuerpo (o en este caso, un equipo) experimente una aceleración y gane la competencia, debe existir una fuerza neta externa que actúe sobre él en la dirección del movimiento deseado. Si consideramos a un solo equipo como nuestro sistema de estudio, las únicas fuerzas horizontales que actúan sobre ellos son: la tensión de la soga (que los jala hacia adelante, hacia el equipo contrario) y la fuerza de fricción del suelo (que los empuja hacia atrás). 

\subsubsection*{Interacción con el Suelo (Tercera Ley de Newton)}
Para generar esa fuerza de fricción estática hacia atrás, el equipo debe interactuar con el suelo. Los competidores aplican una fuerza muscular para empujar el suelo hacia adelante (fuerza de acción). Como consecuencia directa de la Tercera Ley de Newton, el suelo responde aplicando una fuerza de igual magnitud pero en dirección opuesta sobre los zapatos de los competidores, es decir, hacia atrás (fuerza de reacción). El equipo que sea capaz de empujar el suelo con mayor intensidad recibirá, a cambio, una fuerza de reacción externa mayor por parte del suelo. Si esta fuerza de reacción externa supera en magnitud a la fuerza de tensión de la soga, el equipo experimentará una fuerza neta y se acelerará hacia atrás, ganando la competencia.

\subsubsection*{Representación Gráfica (Código Python)}
Para visualizar mejor por qué las fuerzas externas definen al ganador, el siguiente código en Python genera un diagrama de fuerzas de ambos equipos. Se utiliza para demostrar gráficamente que la tensión es una fuerza interna que se anula en el sistema global, mientras que las fuerzas de fricción dictan la aceleración del centro de masa.

\begin{lstlisting}[language=Python]
import matplotlib.pyplot as plt

fig, ax = plt.subplots(figsize=(9, 4))
ax.plot([2.5, 6.5], [1, 1], color='black', linewidth=3)

ax.add_patch(plt.Rectangle((1.5, 0.5), 1, 1, color='blue', alpha=0.5))
ax.add_patch(plt.Rectangle((6.5, 0.5), 1, 1, color='red', alpha=0.5))

ax.annotate('', xy=(3.5, 1), xytext=(2.5, 1), arrowprops=dict(facecolor='black', width=1, headwidth=6))
ax.annotate('', xy=(5.5, 1), xytext=(6.5, 1), arrowprops=dict(facecolor='black', width=1, headwidth=6))
ax.text(4.5, 1.2, 'Tension (Fuerzas Internas Iguales)', ha='center')

ax.annotate('', xy=(0.5, 0.75), xytext=(1.5, 0.75), arrowprops=dict(facecolor='green', width=3, headwidth=10))
ax.text(1, 0.4, 'Reaccion del Suelo (Mayor)', ha='center', color='green', fontsize=9, fontweight='bold')

ax.annotate('', xy=(7.5, 0.75), xytext=(7.5, 0.75), arrowprops=dict(facecolor='orange', width=1, headwidth=5))
ax.annotate('', xy=(7.5, 0.75), xytext=(6.5, 0.75), arrowprops=dict(facecolor='orange', width=1, headwidth=5))
ax.text(7, 0.4, 'Reaccion del Suelo (Menor)', ha='center', color='orange', fontsize=9)

ax.text(2, 1.8, 'Equipo Ganador', ha='center', fontweight='bold', color='blue')
ax.text(7, 1.8, 'Equipo Perdedor', ha='center', fontweight='bold', color='red')

ax.set_xlim(0, 9)
ax.set_ylim(0, 2.5)
ax.axis('off')
plt.title('Diagrama de Fuerzas en una Guerra de Jalar la Soga', pad=20)
plt.show()
\end{lstlisting}

\subsubsection*{Conclusión}
Opción correcta b.


\end{document}
